% =====================================================================
% EXACT SU(N) FLUX-BAND SPECTRA --- standalone
%
% SELF-CONTAINED EDITION.  Generated by paper/master/assemble_selfcontained.py
% from the modular sources in paper/master/.  Do not edit this file: edit
% the section files and re-assemble, or the two will disagree.
%
% Needs no other file.  Build with:
%     tectonic workhouse_band_paper_selfcontained.tex
% or  pdflatex workhouse_band_paper_selfcontained.tex   (twice, for cross-references)
%
% Standard packages only, no shell escape, no images, no .bib.
% =====================================================================
% ===================================================================
% Part I, extracted as a standalone paper.
%
%     tectonic paper/master/workhouse_band_paper.tex
%
% This driver compiles exactly the Part I source files of the master
% edition with no other changes, so the two editions cannot drift.  The
% \onlypartone switch suppresses forward references into Parts II and
% III, which are the only text that differs between the two builds.
% ===================================================================
\documentclass[11pt,letterpaper]{article}
\newcommand{\onlypartone}{1}
% ---- begin preamble.tex --------------------------------------------
% Shared preamble for both drivers: the monograph and the extractable Part I.
% Conventions follow paper/workhouse_publication_edition_rev5_2026-08-30.tex
% so that material can move between editions without renaming symbols.

\usepackage[T1]{fontenc}
\usepackage[utf8]{inputenc}
\usepackage{lmodern}
\usepackage[margin=1in,headheight=14pt]{geometry}
\usepackage{amsmath,amssymb,amsthm,mathtools}
\usepackage{booktabs,array,longtable,enumitem}
\usepackage{xcolor,microtype,xurl}
\usepackage{fancyhdr}
\usepackage[hidelinks]{hyperref}

\hypersetup{colorlinks=true,
  linkcolor=blue!45!black, citecolor=blue!45!black, urlcolor=blue!45!black}

\setlength{\emergencystretch}{3em}
\setlength{\parskip}{2pt}
\setlist{itemsep=2pt,topsep=4pt}
\numberwithin{equation}{section}

\theoremstyle{plain}
\newtheorem{theorem}{Theorem}[section]
\newtheorem{proposition}[theorem]{Proposition}
\newtheorem{lemma}[theorem]{Lemma}
\newtheorem{corollary}[theorem]{Corollary}
\newtheorem{conjecture}[theorem]{Conjecture}
\theoremstyle{definition}
\newtheorem{definition}[theorem]{Definition}
\newtheorem{problem}{Problem}
\theoremstyle{remark}
\newtheorem{remark}[theorem]{Remark}
\newtheorem{scopenote}[theorem]{Scope}

% ---------------------------------------------------------------- notation
\newcommand{\R}{\mathbb{R}}
\newcommand{\C}{\mathbb{C}}
\newcommand{\Q}{\mathbb{Q}}
\newcommand{\Z}{\mathbb{Z}}
\newcommand{\N}{\mathbb{N}}
\newcommand{\SU}{\mathrm{SU}}
\newcommand{\Tr}{\operatorname{Tr}}
\newcommand{\spec}{\operatorname{spec}}
\newcommand{\Ric}{\operatorname{Ric}}
\newcommand{\im}{\operatorname{im}}
\newcommand{\rank}{\operatorname{rank}}
\newcommand{\norm}[1]{\left\lVert #1 \right\rVert}
\newcommand{\abs}[1]{\left\lvert #1 \right\rvert}
\newcommand{\ip}[2]{\left\langle #1, #2 \right\rangle}
\newcommand{\Hil}{\mathcal{H}}
\newcommand{\Alg}{\mathcal{A}}
\newcommand{\gp}{\mathfrak{g}}
\newcommand{\Hph}{\Hil_{\mathrm{phys}}}

% ------------------------------------------------- verification-tier badges
% The badge records how a statement is machine-certified in the repository.
% It is INDEPENDENT of mathematical status: a T3 badge on a proven theorem
% means the proof is analytic and no machine check covers the whole
% statement.  Conflating the two is the error this notation exists to stop.
\definecolor{tierzero}{RGB}{20,90,40}
\definecolor{tierone}{RGB}{25,70,140}
\definecolor{tiertwo}{RGB}{130,90,10}
\definecolor{tierthree}{RGB}{100,100,105}

\newcommand{\tierbadge}[2]{%
  \textcolor{#1}{\footnotesize\textsf{[#2]}}}
\newcommand{\Tzero}{\tierbadge{tierzero}{T0 Lean}}
\newcommand{\Tone}{\tierbadge{tierone}{T1 exact}}
\newcommand{\Ttwo}{\tierbadge{tiertwo}{T2 numeric}}
\newcommand{\Tthree}{\tierbadge{tierthree}{T3 analytic}}

% Graph identity of a statement, so a reader or an agent can run
% `workhouse why <id>` on anything asserted here.
\newcommand{\gid}[1]{\par\nobreak\smallskip\noindent
  {\footnotesize\textsf{Graph id:} \nolinkurl{#1}}\par\smallskip}
\newcommand{\gidc}[2]{\par\nobreak\smallskip\noindent
  {\footnotesize\textsf{Graph id:} \nolinkurl{#1}\quad
   \textsf{Reproduce:} \texttt{#2}}\par\smallskip}
\newcommand{\srcpath}[1]{\nolinkurl{#1}}

% An explicit, visible statement of what a theorem does NOT say.  Used
% wherever the corpus records a scope limit, which is most places.
\newcommand{\notasserted}[1]{\par\smallskip\noindent
  {\small\textbf{Does not assert.} #1}\par\smallskip}

% The obligation register's problem headers.
\newcommand{\obligation}[2]{\par\medskip\noindent
  \textbf{#1}\ \ \textit{#2}\par\smallskip}
% ---- end preamble.tex ----------------------------------------------

\pagestyle{fancy}\fancyhf{}
\fancyhead[L]{\small Exact SU(N) flux-band spectra}
\fancyhead[R]{\small 12 September 2026}
\fancyfoot[C]{\thepage}

% ---- begin gen_macros.tex ------------------------------------------
% Generated by paper/master/generate_apparatus.py.  Do not edit.
% Every count the prose quotes is defined here, from the ledgers, so
% the paper cannot state a number the repository has moved past.
\newcommand{\LitPapers}{108}
\newcommand{\LitEdges}{152}
\newcommand{\LitVerified}{54}
\newcommand{\LitObtained}{34}
\newcommand{\RouteTotal}{111}
\newcommand{\RouteDone}{64}
\newcommand{\RouteLive}{20}
\newcommand{\RouteUntried}{12}
\newcommand{\RouteDead}{15}
\newcommand{\GapOpen}{20}
\newcommand{\GapDischarged}{5}
\newcommand{\ResTotal}{143}
\newcommand{\LeanTotal}{439}
\newcommand{\LeanSorry}{0}
\newcommand{\ContradictionTotal}{22}
\newcommand{\ContraFalsified}{4}
\newcommand{\ContraSuperseded}{3}
\newcommand{\ContraResolved}{15}
\newcommand{\ResConditional}{1}
\newcommand{\ResDisputed}{1}
\newcommand{\ResFalsified}{1}
\newcommand{\ResProven}{140}
% ---- end gen_macros.tex --------------------------------------------
% ---- begin gen_srcmacros.tex ---------------------------------------
% Generated by paper/master/generate_sources.py.  Do not edit.
% Counts for the source-provenance sections.
\newcommand{\CiteDocs}{40}
\newcommand{\DerivStatements}{332}
\newcommand{\DerivOpen}{63}
\newcommand{\LeanModules}{22}
\newcommand{\LibraryPapers}{260}
\newcommand{\LibraryVolOne}{99}
\newcommand{\LibraryVolTwo}{161}
\newcommand{\ExtractedPages}{777}
\newcommand{\ExtractedWorks}{28}
\newcommand{\ExtractedTopics}{7}
\newcommand{\SymbolCount}{28}
\newcommand{\AliasCount}{248}
\newcommand{\NoteArchives}{17}
\newcommand{\NoteReviews}{313}
\newcommand{\GraphEdges}{32503}
\newcommand{\GraphEdgeKinds}{0}
% ---- end gen_srcmacros.tex -----------------------------------------

\title{\bfseries Exact Flux-Band Spectra of Hamiltonian \texorpdfstring{$\SU(N)$}{SU(N)}
Lattice Gauge Theory\\[4pt]
\large The homological carrier, its all-rank coefficients,
and the grading of the strong-coupling series}
\author{Alexander Smith\\\small Independent Researcher}
\date{12 September 2026}

\begin{document}
\maketitle
% ---- begin 00_abstract_band.tex ------------------------------------
\begin{abstract}
\noindent
We study the strong-coupling spectrum of the Kogut--Susskind
Hamiltonian for pure $\SU(N)$ gauge theory on a cubic spatial lattice,
organized around a homological object fixed before any perturbation
theory: the kernel of the plaquette boundary map,
$Z_2 = \ker\partial_2$.

Projection to the fundamental charge-odd one-plaquette manifold gives a
three-orbital Bloch Hamiltonian whose second-order hopping is not an
engineered flat-band model. Non-abelian group integration generates the
oriented edge--face boundary Gram operator $BB^{\dagger}$ with
$B = \partial_2^{\dagger}$, and the amplitude in front of it is the
exact rational function
\[
  t_N = \frac{2N(N^2-4)}{(N^2-1)(2N^2-1)(4N^2-9)} > 0
  \qquad (N \ge 3),
\]
so that colour dynamics fixes the coefficient and cellular geometry
fixes the kernel and the spectrum. The four channel weights producing
$t_N$ are order-two Weingarten values, not an isotropy assumption; the
vanishing at $N = 2$ is an identity between the parallel and
antiparallel channel sums, not a smallness. The projected spectrum
carries one momentum-independent level and a doubly degenerate branch
shifted by $t_N u^2 q(k)$, and the flat carrier has exact multiplicity
$L^3 + 2$ for $L \ge 3$.

At fourth order the effective operator splits into Hodge terms plus one
geometric term coupling the carrier to its complement, with support
$\{I, U, S, S^2, R\}$. Its off-axis coefficient was disputed between
two exact rationals and is decided here by derivation: three
independent implementations and the original pipeline's own recovered
intermediate ledger agree on every cluster but one, and the exception is
diagnosed as an enumeration shortfall of sixteen missing insertion
orderings, giving
$C_{\mathrm{shp}} = C_{\mathrm{historical}} + 25/1024$ exactly.

Every coefficient is computed as a rational function of the rank rather
than interpolated through ranks: we run the exact engine over $\Q(N)$,
which removes the degree bound that rational interpolation would
otherwise require. This yields two grading theorems. At even orders the
one-plaquette $C$-parity splitting obeys
$S_m(N) = c_m N^{-(2m-1)}(1 + O(N^{-2}))$ for $m \le 10$, with
$c_2 = -4$, $c_4 = -32$, $c_6 = -1748/3$, $c_8 = -123332/9$,
$c_{10} = -49593808/135$, and the scaling is a cancellation theorem:
four intermediate channels each of order $N^{-(m-1)}$ have leading
coefficients in the ratio $(-4,+2,+1,+1)$, summing to zero. At odd
orders the coefficients vanish identically for even $N$; for odd $N$ the
first survivor is order $N-2$, a single closed-form vertex.

At sixth order we give the complete eighteen-word folded operator and
prove, by exhibiting an explicit point of the Laurent ring where $q$
vanishes and the numerator does not, that the sixth-order band
coefficient cannot lie in the fourth-order shape span --- whatever the
sixth-order dynamics turns out to be.

These are finite-order, finite-volume statements about coefficients of a
formal expansion, established exactly. They are not a continuum
glueball theorem and not a mass-gap theorem. Every statement carries its
machine-verification tier and the command that re-establishes it.
\end{abstract}
% ---- end 00_abstract_band.tex --------------------------------------
\tableofcontents
\bigskip

% ---- begin 10_setting.tex ------------------------------------------
\section{Setting and conventions}
\label{sec:setting}

\subsection{The Hamiltonian}

Work on the cubic lattice $\Z^3$ with gauge group $G = \SU(N)$, one
group element $U_\ell$ per oriented link. The Kogut--Susskind
Hamiltonian is
\begin{equation}
  H \;=\; H_0 + V,
  \qquad
  H_0 = \tfrac12 \sum_{\ell} E_\ell^2,
  \qquad
  V = -\,u\, W,
  \label{eq:KS}
\end{equation}
where $E_\ell^2$ is the quadratic Casimir of the left-invariant electric
field on link $\ell$, and $W$ is the sum of plaquette characters in both
orientations,
\begin{equation}
  W = \sum_{p} \big(\chi_p + \bar\chi_p\big),
  \qquad
  \chi_p = \Tr \prod_{\ell \in \partial p} U_\ell .
\end{equation}

Strong-coupling perturbation theory is the expansion in $u$ around the
electric vacuum, organized as degenerate perturbation theory --- a des
Cloizeaux effective operator --- on the sector of interest.

\subsection{The coupling, and one erratum that matters}

The canonical coupling variable throughout is
\begin{equation}
  u \;=\; \frac{\beta_N}{2N},
  \label{eq:coupling}
\end{equation}
whose $\SU(3)$ specialization is $u = \beta_3/6$.

A warning, because it has corrupted transcriptions of this material
before. An archived line reading $Y = 2\beta/3 = 4u$ is a
definition-label erratum, not a change of variable. The coefficients
printed alongside it were \emph{already} in the canonical coordinate
\eqref{eq:coupling}. Applying a $4^r$ rescaling at order $r$ to
``correct'' them corrupts every order, and the repository carries an
explicit check that the printed towers reproduce their canonical values
verbatim.

For rank-uniform statements the natural variable is the planar one,
\begin{equation}
  \tau \;=\; \frac{2u}{N^2} \;=\; \frac{\beta_N}{N^3},
  \label{eq:tau}
\end{equation}
in which the grading theorems of \S\ref{sec:grading} take their simplest
form: every even order of the one-plaquette splitting then carries the
same single power of $N$.

\subsection{Two projections that must not be conflated}

Two distinct projections appear throughout, and several contradictions
in this program's own register were naming collisions between them.

The \emph{band} projection is onto the degenerate strong-coupling
subspace carrying the excitation of interest --- for the flux band, the
cellular subspace of \S\ref{sec:carrier}. Coefficients anchored to it
are band-kernel quantities.

The \emph{physical} projection is onto the gauge-invariant,
vacuum-subtracted spectrum at a definite momentum. Coefficients anchored
to it are physical quantities.

These are related by translation-local scalar shifts, and such a shift
changes neither the centered operator, nor its eigenvectors, nor the
bandwidth. It does change the scalar. Consequently the band-kernel
anchor and the physical $\Gamma$-point coefficient are two differently
anchored coordinates, not two competing estimates of one quantity, and
they are written $q^{(4)}_{\mathrm{band}}$ and $m^{(4)}_{\Gamma}$
respectively. Writing either as ``$m_4$'' regenerates a contradiction
that does not exist, and the repository's search tooling refuses the
name for that reason.

\subsection{Exactness}

Every coefficient displayed in Part~\ref{part:exact} is an exact
rational, or an exact rational function of $N$. Where a quantity is
known here only numerically it carries an explicit numerical marker in
the source registry and is never displayed as though exact. The single
most dangerous defect available in this material is a float that reads
as exact, and the boundary between the two is machine-guarded rather
than maintained by care.
% ---- end 10_setting.tex --------------------------------------------
% ---- begin 11_carrier.tex ------------------------------------------
\section{The homological carrier}
\label{sec:carrier}

\subsection{Where the flat band comes from}

Project the Kogut--Susskind Hamiltonian \eqref{eq:KS} onto the
fundamental, charge-odd, one-plaquette manifold. In three spatial
dimensions a fundamental plaquette excitation carries one of three face
orientations, and at second order neighbouring faces communicate
through a shared link. The projected problem is therefore a
three-orbital Bloch Hamiltonian, and it looks at first like a
tight-binding model.

It is not an engineered one. The signs of the matrix elements are fixed
by the orientation of the plaquette boundary and by charge conjugation,
and non-abelian group integration generates, exactly, the oriented
edge--face boundary Gram operator
\begin{equation}
  B(k)B(k)^{\dagger},
  \qquad B = \partial_2^{\dagger},
\end{equation}
with the coefficient of \S\ref{sec:second} in front of it. The
structure of the result is a clean separation:
\begin{equation}
  \boxed{\ \text{colour dynamics} \longrightarrow t_N,
  \qquad
  \text{cellular geometry} \longrightarrow BB^{\dagger}. \ }
  \label{eq:separation}
\end{equation}
Representation theory fixes the amplitude; the cubical chain complex
fixes the kernel and the spectrum. Neither borrows from the other.

\subsection{The projected spectrum}

\begin{theorem}[Second-order projected spectrum]
\label{thm:projspec}
The projected second-order spectrum consists of one
momentum-independent level together with a doubly degenerate branch
shifted by $t_N u^2 q(k)$, where
\begin{equation}
  q(k) \;=\; 4\sum_{j=1}^{3} \sin^2\!\big(k_j/2\big).
\end{equation}
\end{theorem}

The flat level is exact and it is exact for a reason: the flat carrier
is $\ker \partial_2$.

\begin{theorem}[The carrier and its multiplicity]
\label{thm:multiplicity}
Elementary cube boundaries are compact localized states of the carrier;
three wrapping sheets complete it on the torus. For $L \ge 3$ the exact
multiplicity is
\begin{equation}
  (L^3 - 1) + 3 \;=\; L^3 + 2 .
\end{equation}
\end{theorem}
\gid{RESULT:UNIVERSAL\_CELLULAR\_HODGE\_SPECTRUM}

The two contributions are structurally different and the distinction
matters later. The $L^3 - 1$ cube boundaries are local and finitely
supported. The three sheets are not: they wrap the torus, and they are
the entire zero-momentum content of the carrier. \S\ref{sec:wilson}
returns to the consequence --- that no finitely supported exactly-flat
local operator can see the $\Gamma$ point at all.

\subsection{Third order, and the coefficients that survive}

For $\SU(3)$, exact third-order perturbation theory stays inside the
same boundary-factorized algebra. It does not generate a new operator
structure; it renormalizes the coefficient in front of the same
$BB^{\dagger}$:
\begin{equation}
  H_{\mathrm{eff},-}
  \;=\; E_{\mathrm{flat}}(u)\, I
  \;+\; \left(\frac{5}{612}u^2 + \frac{1975}{124848}u^3\right) BB^{\dagger}
  \;+\; O(u^4),
  \label{eq:thirdorder}
\end{equation}
with
\begin{equation}
  E_{\mathrm{flat}}(u)
  \;=\; \frac{8}{3} + u + \frac{11}{306}u^2
        - \frac{109151}{249696}u^3 .
\end{equation}

That \eqref{eq:thirdorder} remains boundary-factorized through third
order is the non-trivial part. It is what makes the fourth order the
first place where the carrier can couple to its complement, and
\S\ref{sec:fourth} is about exactly that coupling.

\subsection{What is and is not novel here}

This program does not claim that incidence flatness, cellular compact
localized states, or homological counting of zero modes are new. All
three have a substantial prior literature --- line-graph constructions,
compact localized states, the necessity of non-contractible states on
periodic systems, rectangular factorizations and index arguments, and
closely related closed-surface structures in two-form spin liquids ---
and a recent independent construction treats face-graph and higher-cell
flat bands directly.

What is claimed is that the flatness here is \emph{dynamical}: it is
produced by exact strong-coupling matrix elements of non-abelian
$\SU(N)$ Hamiltonian lattice gauge theory, with the amplitude fixed by
group representation theory and the kernel fixed by the cubical chain
complex, as in \eqref{eq:separation}. The flat band is not put in; it
comes out.

\begin{scopenote}
Theorem~\ref{thm:projspec} and \eqref{eq:thirdorder} are finite-volume,
finite-order statements about a projected flux sector. They are not a
continuum glueball theorem and not a mass-gap theorem, and the source
manuscript says so in its own abstract.
\end{scopenote}
% ---- end 11_carrier.tex --------------------------------------------
% ---- begin 12_second_order.tex -------------------------------------
\section{Second order at every rank}
\label{sec:second}

This section is the algebraic base of the whole program. Everything in
it is proof-checked, and its conclusion --- that the band exists for
$N \ge 3$ and not for $N = 2$, with an exact amplitude --- follows from
representation theory with no input from any expansion.

\subsection{The shared-link weights are Weingarten values}

Two plaquettes contributing at second order either share a link or do
not. The amplitude for the shared-link case was for a long time taken
in this program with an isotropy premise: that the four contributing
representation channels carry weights in a fixed ratio. That premise
is not needed. It is a consequence of the order-two Weingarten values.

\begin{proposition}[The shared-link weights, from representation theory alone]
\label{prop:weingarten}
Let two plaquettes share exactly one link. Then:
\begin{enumerate}
  \item The six non-shared links collapse. Each plaquette contributes a
    product of three independent Haar links, and a product of
    independent Haar matrices is Haar. So the two-plaquette amplitude
    reduces to $\Tr(A U)\,\Tr(B U^{\pm 1})$ with $A$ and $B$ independent
    Haar, and integrating them out leaves a pure degree-$(2,2)$ moment of
    the shared link $U$.
  \item Two such moments settle both families. With
    $M_{\mathrm{direct}} = N^2$ and $M_{\mathrm{cross}} = N$, the like
    family splits as $(N+1)/(2N)$ and $(N-1)/(2N)$; and the singlet
    component of $U_{ij}\overline{U_{lk}}$ is
    $\delta_{il}\delta_{jk}/N$, of squared norm $1$, giving $1/N^2$.
\end{enumerate}
All four weights are exactly $d_R/N^2$, where $d_R$ is the dimension of
the channel.
\end{proposition}
\gidc{the shared-link weights are Weingarten, not an isotropy assumption}{workhouse verify --only 'the shared-link weights are Weingarten, not an isotropy assumption'}

\noindent\Tone{}. The relevant Weingarten pair is the inverse of the
$S_2$ Gram matrix and the index sums are explicit, so the check imports
nothing from the corpus. Registered as the discharge of the gap that
had carried this premise as the one unproved input of the second-order
chain; it was opened and closed on the same day, 28~August 2026, once
the question was asked in the right form.

\subsection{The four channels in closed form}

Write $C_F = (N^2-1)/(2N)$ for the fundamental Casimir. Each channel
$R$ contributes a resolvent weight $-(d_R/N^2)/(C_F + \tfrac12 \Delta C_2(R))$,
and all four evaluate in closed form.

\begin{theorem}[Per-channel resolvent weights]
\label{thm:channels}
\begin{align}
  w_{\mathbf 1}   &= -\frac{1/N^2}{C_F}
    &&= -\frac{2}{N(N^2-1)},
  \\[3pt]
  w_{\mathrm{Adj}} &= -\frac{(N^2-1)/N^2}{C_F + N/2}
    &&= -\frac{2(N^2-1)}{N(2N^2-1)},
  \\[3pt]
  w_{\mathrm{Sym}} &= -\frac{N(N+1)/2N^2}
                        {C_F + \frac{(N-1)(N+2)}{2N}}
    &&= -\frac{N+1}{(N-1)(2N+3)},
  \\[3pt]
  w_{\Lambda} &= -\frac{N(N-1)/2N^2}
                    {C_F + \frac{(N+1)(N-2)}{2N}}
    &&= -\frac{N-1}{(N+1)(2N-3)} .
\end{align}
\end{theorem}
\gid{LEAN:resolventWeight\_singlet, LEAN:resolventWeight\_adjoint, LEAN:resolventWeight\_sym, LEAN:resolventWeight\_antisym}

\noindent\Tzero{} --- all four, at
\srcpath{lean/Workhouse/Basic.lean:72,80,89,104}. Reproduce with
\texttt{make lean}.

\subsection{The two family sums}

The channels group into a \emph{mixed} family $\{\mathbf 1, \mathrm{Adj}\}$
and a \emph{like} family $\{\Lambda^2 F, \mathrm{Sym}^2 F\}$, and these
are the antiparallel and parallel channel sums.

\begin{theorem}[Family sums]
\label{thm:families}
\begin{equation}
  A_N \;=\; w_{\mathbf 1} + w_{\mathrm{Adj}}
    \;=\; -\,\frac{2N^3}{(N^2-1)(2N^2-1)},
  \qquad
  B_N \;=\; w_{\mathrm{Sym}} + w_{\Lambda}
    \;=\; -\,\frac{4N(N^2-2)}{(N^2-1)(4N^2-9)} .
\end{equation}
\end{theorem}
\gid{LEAN:mixed\_family\_sum, LEAN:like\_family\_sum}

\noindent\Tzero{}, at \srcpath{lean/Workhouse/Basic.lean:118,129}. The
mixed sum is obtained from the dimension and Casimir table rather than
by transcription from any source document, which is the point: the
family sums are re-derived, not quoted.

\subsection{The hopping amplitude}

\begin{theorem}[The all-rank second-order hopping amplitude]
\label{thm:tN}
\begin{equation}
  t_N \;=\; B_N - A_N
  \;=\; \frac{2N\,(N^2-4)}{(N^2-1)(2N^2-1)(4N^2-9)} .
  \label{eq:tN}
\end{equation}
It is strictly positive for every $N \ge 3$, and
\begin{equation}
  t_2 = 0, \qquad t_3 = \tfrac{5}{612}.
\end{equation}
\end{theorem}
\gidc{CONST:t\_N}{workhouse why t\_N}

\noindent\Tzero{} for the rank law and for both special values
(\srcpath{lean/Workhouse/Basic.lean}: \texttt{rank\_law\_numerator},
\texttt{hopping\_two}, \texttt{hopping\_three}). Clearing denominators,
the channel difference has numerator exactly $2N(N^2-4)$.

Two consequences deserve to be stated separately, because both are
structural rather than numerical.

\begin{corollary}[$\SU(2)$ has no band at this order]
\label{cor:su2}
The factor $N^2 - 4$ in \eqref{eq:tN} vanishes at $N = 2$. The exclusion
of $\SU(2)$ is therefore not a smallness but an identity: the parallel
and antiparallel channel sums coincide exactly at rank two.
\end{corollary}
\gid{LEAN:hopping\_two}

\begin{corollary}[The planar approach is from below]
\label{cor:deficit}
With $D = (N^2-1)(2N^2-1)(4N^2-9)$,
\begin{equation}
  D - 4N^3 \cdot 2N(N^2-4) \;=\; 2N^4 + 31N^2 - 9 \;>\; 0
  \qquad (N \ge 1),
\end{equation}
so $4N^3 t_N = 1 - (2N^4+31N^2-9)/D < 1$. Hence
$t_N < 1/(4N^3)$ for every $N$, and $t_N \sim 1/(4N^3)$: the
$N^{-3}$ cancellation is exact and what it leaves is a strictly positive
deficit.
\end{corollary}
\gid{LEAN:hopping\_deficit\_numerator}

\noindent\Tzero{}, at \srcpath{lean/Workhouse/Basic.lean:48}.
Corollary~\ref{cor:deficit} is the first appearance of the rank grading
that \S\ref{sec:grading} develops: the leading power of $N$ is fixed by
a cancellation, and the surviving term is a polynomial one can write
down.

\subsection{Independent corroboration of the channel weights}

The weights of Proposition~\ref{prop:weingarten} have an external
provenance, from work written without knowledge of this program.

An efficient-truncation study of $\SU(N_c)$ lattice gauge theory for
quantum simulation~\cite{CBB_2026} gives, in its Appendix~D, the exact
finite-rank plaquette matrix element
\begin{equation}
  \abs{M_\rho} = \sqrt{\frac{\dim\rho}{\dim A \,\dim R}} .
\end{equation}
Taking both factors fundamental gives $\abs{M_\rho}^2 = d_\rho/N^2$,
which is \emph{identically} the numerator of the weight formula used
here. The four shared-link channel weights, and with them $A_N$, $B_N$
and $t_N = B_N - A_N$, therefore rest on a matrix element an
independent group wrote down for an unrelated purpose. That truncation
is also the first published one retaining all four shared-link
irreducible representations.

The agreement is evidence rather than bookkeeping precisely because
neither side knew of the other. It does not make $t_N$ published: the
rank law \eqref{eq:tN}, the vanishing at $N = 2$, and the deficit of
Corollary~\ref{cor:deficit} are assembled here.

A caution attaches in the other direction. The leading-large-$N_c$
qutrit plaquette Hamiltonian~\cite{CB_2024} explicitly decouples the
charge-odd combination, so the sector in which $t_N$ lives is inert in
that model --- effectively $t_3 = 0$ at leading order in large $N_c$,
against the finite-$N$ channel-complete value $5/612$. That paper is
recorded in this program's index as not yet obtained, and the statement
here is taken from imported bridge documents rather than checked
against the paper itself.

\subsection{What this section does not establish}

\notasserted{Convergence of the strong-coupling series; anything about
the spectrum at finite $u$ beyond second order; and nothing about
$\SU(2)$ at higher orders, where Corollary~\ref{cor:su2} does not
apply. Theorem~\ref{thm:tN} is a statement about one coefficient of a
formal expansion, established exactly.}
% ---- end 12_second_order.tex ---------------------------------------
% ---- begin 13_fourth_order.tex -------------------------------------
\section{Fourth order, and how a disputed coefficient was decided}
\label{sec:fourth}

Fourth order is where the carrier first couples to its complement, and
it is where this program spent the largest single block of its effort.
The mathematics is one theorem and one number. The number was disputed
for months between two irreconcilable values, and the way it was settled
is, in this author's view, the more instructive half of the story.

\subsection{The structure}

The fourth-order effective operator splits into Hodge terms and one
geometric term coupling the carrier to its complement. Restricted to
range-one operator structures, the shape content collapses: each such
structure contributes a fixed shape coefficient, and every one of them
gives $B = D = 0$ identically. The tier collapse is simply what
range-one looks like.

\begin{theorem}[Fourth-order support, at $\SU(3)$]
\label{thm:tiercollapse}
The recorded fourth-order kernel has support
$\{I, U, S, S^2, R\}$, and the shape data closes in that span. Solved
over the whole Brillouin zone in exact rationals rather than fitted, the
shape table is determined by three signed amplitudes:
\begin{equation}
  C_{\mathrm{shp}} \;=\; -\tfrac{5}{96} - u - \tfrac12(\rho + \pi),
  \qquad
  u_2 = 2u, \qquad \nu = -\big(\tfrac{5}{48} + 4u\big),
  \label{eq:kernelform}
\end{equation}
where $u$ is the two-hop chain cumulant.
\end{theorem}
\gid{RESULT:TIER\_COLLAPSE\_ACTUAL\_H4\_SUPPORT, RESULT:TIER\_COLLAPSE\_IS\_R\_DEGREE}

\begin{scopenote}
``The recorded kernel'' means the $189$-record historical kernel, which
is $N = 3$. The operator splitting is an $\SU(3)$ statement; it is not
established at general rank, and framing it as an $\SU(N)$ identity
would be an overreach. The rank-uniform statements of this program are
the coefficient laws of \S\ref{sec:second} and \S\ref{sec:grading}, not
this splitting.
\end{scopenote}

Clearing the $1/q$ makes the shape ansatz a linear system over Laurent
polynomials in $e^{ik_j}$, which is why the whole zone can be solved
exactly instead of sampled.

There is a second, equivalent presentation. In the Hodge form the
two-hop weight is absorbed, $\tilde\pi = \pi + 2u$, and
\begin{equation}
  C_{\mathrm{shp}} \;=\; -\tfrac{5}{96} - \tfrac12\big(\rho + \tilde\pi\big),
  \label{eq:hodgeform}
\end{equation}
so that $u$ drops out of $C$ entirely and exactly two shared-link
amplitudes remain. Equations \eqref{eq:kernelform} and
\eqref{eq:hodgeform} are the same statement. They must not be
combined --- writing $u$ and $\tilde\pi$ in the same expression
double-counts the two-hop weight, and that error has been made.

\subsection{The dispute}

Two independently recorded kernels --- a historical pipeline from June
2026 and a later cold dump --- gave different values for the off-axis
shape coefficient $C_{\mathrm{shp}}$. Both were exact rationals. Neither
could be promoted, because the repository's governing rule is that a
dispute closes by derivation or not at all: code may not pick a side,
average them, or prefer whichever rational looks more authoritative.

Comparing the two kernels record by record localized the disagreement
rather than resolving it. Both have the same six weight orbits and the
same normalized shape table; both satisfy the same two relations above.
What differs is the skeleton scale and the \emph{signs} of two
amplitudes. One scale factor explains $144$ of the $189$ records to
$2\times10^{-12}$: the rotation bulk is structurally shared, and the
dispute is three amplitudes, not two kernels.

\subsection{How it was settled}

By building a third implementation that shares no primitive with either,
and by recovering the historical pipeline's own intermediate ledger.

The third implementation is a Wilson-loop calculus written from scratch:
the free Hamiltonian by the Fierz identity as a rewiring of index ports
--- a like pair swaps wires, an unlike pair is cut and crossed, which is
unitarity; Haar integrals by the $U(N)$ Weingarten function as the Gram
pseudoinverse on $S_n$, with $\det U = 1$ inserted on virtual slots for
the determinant families; and the resolvent by per-link irreducible
projectors, so that every component is an exact eigenvector of $H_0$
with a known energy. It reproduces the four second-order constants, the
two-hop weight with the correct sign on three chain types, and the
single-contact and shared-link dressings, in seconds --- and then gives
the corner cumulant to the digit.

Separately, a survey of the corpus located the historical pipeline's own
Stage-3I ledger: $4{,}221$ ordered words whose hash matches the input
the kernel copies quote. Its rooted assembly, ported verbatim, reproduces
the pinned $189$ records exactly. Grouping each word's rooted images by
the set of plaquettes it touches then yields the historical kernel's own
connected cumulants --- a reading the original pipeline never performed
on its own data.

\begin{theorem}[Adjudication of the off-axis coefficient]
\label{thm:cshp}
The historical kernel and the independent assembly agree on \emph{every}
cluster of the rotation record --- pair, fourteen chains, two fans, two
corners --- with identical rationals in both $C$-parity sectors, except
one. The exception is the adjacent-face cube completion. The historical
ledger holds $8$ of the $24$ insertion orderings, each with exactly the
weight the multi-loop model assigns it, namely
$(2,1,4,2,8,8,2,4)$ in units of $1/1536$, totalling $31/1536$; the
sixteen it lacks sum to $-25/512$. The opposite-face completion in the
same ledger has all $24$.

Consequently, in the kernel's basis,
\begin{equation}
  C_{\mathrm{shp}}
  \;=\; C_{\mathrm{historical}} + \frac{25}{1024}
  \;=\; -\,\frac{13035490122347}{550663802582400}
  \;=\; -0.0236723\ldots
  \label{eq:cshp}
\end{equation}
Substituted into the $24$ records and solved over the whole zone, this
gives $A = 5/48$ and $B = D = 0$ with no residual.
\end{theorem}
\gidc{C2, ADR:0024}{workhouse why C2}

\noindent\Tone{}, suite ``the third implementation and the historical
ledger'': nine exact checks, three of them registered findings.

The shortfall has a closed form at every rank. The adjacent-face
completion is
\begin{equation}
  c_4^{\,\mathrm{adj}}(N) = \frac{-106}{N(N^2-1)^3},
  \qquad
  c_4^{\,\mathrm{adj}}(3) = -\frac{53}{768},
\end{equation}
in ratio $53/80$ to the opposite-face completion $-160/(N(N^2-1)^3)$;
the $-106$ is the primitive $-88$ plus ten multi-loop orderings at
$-18$.

\subsection{Three things this episode establishes}

\emph{The disagreement was an enumeration shortfall, not a
disagreement about mathematics.} A pipeline dropped sixteen of
twenty-four orderings on one cluster and nowhere else. That is a
diagnosable defect, and it was diagnosed ordering by ordering rather
than adjudicated by preference.

\emph{A basis error can survive validation.} An earlier assembly of the
same amplitude was correct and was read in the conjugate basis, giving
a third value. Every $C$-odd component of the historical ledger is the
negative of the assembly's value in one traversal while every $C$-even
one is identical --- the signature of conjugating a single face. The
assembly was right; its basis was not. The decision record amends it
rather than withdrawing it.

\emph{The right answer arrived by a forbidden route first.} The value
\eqref{eq:cshp} is the one an all-rank continuation formula produces at
$N=3$ --- and that derivation was retracted, because the substitution it
used is not permitted. The retraction stands. That the forbidden route
happened to reach the number the legitimate assembly later confirmed is
not a vindication of the route, and the register keeps both facts.

\begin{scopenote}
Both recorded kernels remain in the registry with their verdicts: the
historical superseded, the cold dump falsified, beside the proven
assembled value. Nothing was deleted. One follow-up remains untried ---
the pair cluster from the third engine, pure-six family included ---
which can sharpen this result but not overturn it.
\end{scopenote}
% ---- end 13_fourth_order.tex ---------------------------------------
% ---- begin 14_rank_grading.tex -------------------------------------
\section{Rank grading: cancellation and determinant families}
\label{sec:grading}

The strong-coupling coefficients of \S\ref{sec:second} and
\S\ref{sec:fourth} are exact rational functions of $N$. That makes a
question available which is invisible at any single rank: how does each
order of the series scale with the rank, and why?

The answer has two halves. At even orders the scaling is fixed by an
exact cancellation among four representation channels. At odd orders
the coefficients vanish outright for even $N$, and for odd $N$ they
first appear at order $N-2$ with a closed-form vertex. Neither statement
can be seen from numerics at $N = 3$.

\subsection{The one-face grading law}

Let $S_m(N)$ be the $C$-parity splitting of the one-plaquette
strong-coupling diagonal at order $m$,
\begin{equation}
  S_m(N) \;=\; (\text{odd})_m - (\text{even})_m \;=\; -2\,K_m[0][1].
\end{equation}

\begin{theorem}[Single-face planar grading through tenth order]
\label{thm:grading}
At even orders $m = 2, 4, 6, 8, 10$ the exact rational function of $N$
satisfies
\begin{equation}
  S_m(N) \;=\; c_m\, N^{-(2m-1)}\big(1 + O(N^{-2})\big),
  \label{eq:grading}
\end{equation}
with
\begin{equation}
  c_2 = -4, \quad c_4 = -32, \quad c_6 = -\tfrac{1748}{3},
  \quad c_8 = -\tfrac{123332}{9}, \quad c_{10} = -\tfrac{49593808}{135}.
\end{equation}
Writing $u = N^2\tau/2$ as in \eqref{eq:tau}, the \emph{leading} part of
the order-$m$ contribution is $(c_m/2^m)\, N \tau^m$. Every even order
therefore carries the same single power of $N$, with
\begin{equation}
  b_2 = -1, \quad b_4 = -2, \quad b_6 = -\tfrac{437}{48},
  \quad b_8 = -\tfrac{30833}{576}, \quad b_{10} = -\tfrac{3099613}{8640},
  \qquad b_m = c_m/2^m .
\end{equation}
\end{theorem}
\gid{RESULT:SINGLE\_FACE\_PLANAR\_GRADING}

\noindent Recorded \texttt{proven} with \texttt{output-certified}
evidence. The exponent and the coefficients are exact limits of exact
rational functions of $N$, not numerical extrapolations. An independent
integer-rank engine reproduces $c_2, c_4, c_6, c_8$ by extrapolation to
within $1.4 \times 10^{-5}$ relative --- a cross-check by a different
method, not the source of the values.

The word \emph{leading} in Theorem~\ref{thm:grading} is doing work and
should not be dropped. Each $S_m$ is an exact rational function of $N$
with relative corrections in $N^{-2}$, so the reduction to $\tau$ is a
planar-limit statement, not an identity at finite rank. Explicitly, at
second order
\begin{equation}
  S_2(N) = -4N^{-3} - 4N^{-7},
  \qquad\text{so}\qquad
  \frac{S_2\,u^2}{N} = -\tau^2 - \frac{\tau^2}{N^4}.
\end{equation}
The single-face band per unit $N$ is therefore a function of $\tau$
alone \emph{as} $N \to \infty$ at fixed $\tau$, and not at finite $N$.
A statement of the second form would be false, and an earlier internal
summary of this result made it.

\subsection{The grading is a cancellation theorem}

Equation~\eqref{eq:grading} is a much smaller number than its
constituents, and the reason is exact.

\begin{theorem}[Four-channel cancellation]
\label{thm:cancellation}
Decomposing $S_m$ by the intermediate state of the perturbed vector
gives exactly four contributing channels at every even $m$ from $4$ to
$10$: the singlet $\big((),()\big)$, the adjoint
$\big((1),(1)\big)$, and the two-box states $\big((2),()\big)$ and
$\big((1,1),()\big)$. Each is of order $N^{-(m-1)}$, with leading
coefficients
\begin{equation}
  (-4,\ +2,\ +1,\ +1) \times 2^{\,m/2-2},
\end{equation}
which \emph{sum to zero}. At $m = 2$ only the singlet and the adjoint
contribute, at $-2$ and $+2$, and these also cancel.
\end{theorem}
\gid{RESULT:SINGLE\_FACE\_PLANAR\_GRADING}

So the individual channels are of order $N^{-(m-1)}$ and their leading
behaviour annihilates, leaving \eqref{eq:grading} at $N^{-(2m-1)}$ ---
two further orders down in the $N^{-2}$ expansion. At $m = 4$ this is
the statement that channels of order $N^{-3}$ conspire to produce a band
of order $N^{-7}$.

This is why the grading is stated as a theorem about a mechanism rather
than as an observed power. A power law read off from fitted data at
several ranks would carry no information about \emph{why} the leading
terms are absent, and would give no reason to expect
\eqref{eq:grading} to persist. The cancellation does.

\begin{scopenote}
Theorem~\ref{thm:grading} concerns the \emph{one-plaquette} diagonal
splitting at even orders $m \le 10$. It is not the cluster cumulant of
the multi-face conjecture, which is not discharged by it: the one-face
case at sixth order was already established separately, and the
shared-link pair route carries the multi-face statement. The general-$m$
law remains a conjecture. All statements are about coefficients at fixed
order, not about convergence.
\end{scopenote}

\subsection{Odd orders are determinant families}

\begin{theorem}[Centre parity]
\label{thm:centreparity}
For \emph{even} $N$, every odd-order matrix element of the des Cloizeaux
effective operator of the strong-coupling band vanishes identically ---
on every cluster, in both $C$-parity sectors, at every order.

For \emph{odd} $N$, an odd order $m$ vanishes unless $m \ge N-2$. At
$m = N-2$ the only surviving element is the one-plaquette vertex
\begin{equation}
  H_{N-2}(\bar F, F)
  \;=\; -\,\frac{\big(N/(N+1)\big)^{N-3}}{\big((N-3)!\big)^2},
  \label{eq:centrevertex}
\end{equation}
which is $-1$ at $N = 3$ and $-\tfrac{25}{144}$ at $N = 5$. Hops and
leakages vanish at that order.
\end{theorem}
\gid{RESULT:ODD\_ORDER\_CENTRE\_PARITY}

The mechanism is the centre. The $\SU(N)$ Haar integral of a word
vanishes unless every link's net flux is $0 \bmod N$; on a single
plaquette the four links are private, so the intermediate states are
characters with energy $2C_2(R)$, and the flux condition at odd order
can be met only once enough boxes have accumulated to close modulo $N$.
That threshold is $N-2$, and it is why the first odd order is a
function of the rank rather than a fixed number.

\begin{scopenote}
Statements about coefficients of the strong-coupling series at fixed
order, not about convergence. On a periodic lattice the link-set lemma
requires even extents. The even orders and the overlap theorem are not
addressed here.
\end{scopenote}

\subsection{Why this refuted a claim}

The register records a claim that centre-only circuits are dynamically
dark --- that configurations carrying only centre flux cannot contribute
to the dynamics. Theorem~\ref{thm:centreparity} and its fifth-order
companion refute it: the exact fifth-order determinant correction
$\delta c_{5,\mathrm{det}}$ is nonzero. The claim was falsified by an
exact computation rather than by an estimate, which is the only way a
claim of that shape can be settled.
% ---- end 14_rank_grading.tex ---------------------------------------
% ---- begin 15_symbolic.tex -----------------------------------------
\section{The symbolic-rank engine}
\label{sec:symbolic}

Every closed form in \S\ref{sec:second}--\S\ref{sec:grading} is a
rational function of $N$ computed as such, not interpolated through
ranks. This section explains how, because the method is the most
transferable thing in the program and because getting it wrong is the
default.

\subsection{The problem with sweeping ranks}

The natural procedure is to run an exact engine at $N = 3, 4, 5, \ldots$
and fit a rational function. That is what was done first. An assembled
cluster cumulant $\beta_N$ agreed with a candidate closed form
$P_{17}(N^2)\big/\big(N\,R_{20}(N^2)\big)$ at every rank from $4$ to
$70$: sixty-seven exact agreements between exact rationals.

Sixty-seven exact agreements do not prove the identity. Rational
interpolation is unique only once the degrees are bounded, and without
a degree bound the agreements constrain nothing about $N = 71$. The
review of that work made exactly this objection --- no ``for all $N$''
without a degree bound --- and it was correct.

A degree bound could have been argued: Weingarten denominators, energy
denominators, and a count of vertices together bound the degree, and
the sixty-seven agreements would then prove the identity by uniqueness.
That argument was not made. What was done instead makes it unnecessary,
and is better.

\subsection{Running the engine over \texorpdfstring{$\Q(N)$}{Q(N)}}

Inspect where the rank actually enters the exact engine. It enters in
exactly two places that are not rational arithmetic:
\begin{itemize}
  \item the charge predicate, $\text{charge} \equiv 0 \bmod N$;
  \item the flux-difference predicate, $\abs{\mathrm{diff}} \le N$.
\end{itemize}
Everywhere else the rank appears only through rational expressions in
$C_F = (N^2-1)/(2N)$, in $1/(2N)$, and in $N^{\text{closed}}$.

So make the two predicates overridable and promote everything else to
the rational function field:
\begin{itemize}
  \item scalars become elements of $\Q(N)$, carried as a numerator and
    a monic denominator in exact polynomial arithmetic;
  \item the Weingarten function becomes the symbolic inverse of the
    Gram matrix $N^{\text{cycles}}$ on $S_n$;
  \item the single-link spectra become the $\SU(N)$ quadratic Casimirs
    of the irreducible representations that link's flux content can
    carry.
\end{itemize}

A cumulant then emerges as \emph{one rational function of $N$},
computed rather than fitted. The per-rank engine is otherwise untouched
and its own test suite still holds --- which is the check that the
symbolic run is the same calculation and not a different one.

\begin{theorem}[Closed forms over $\Q(N)$]
\label{thm:symbolic}
The fourth-order cluster cumulants, the two-hop weight in both
$C$-parity sectors, and the assembled $\beta_N$ are exact elements of
$\Q(N)$ produced by the symbolic engine. The identity
$\beta_N = P_{17}(N^2)\big/\big(N\,R_{20}(N^2)\big)$ holds as an
identity of rational functions.
\end{theorem}
\gid{ADR:0029}

\noindent The final polynomial identity is promoted to \Tzero{}, and it
is worth saying exactly which identity that is. The corpus \emph{states}
$\beta_N = P_{17}(N^2)/(N R_{20}(N^2))$ and never derives it. What Lean
checks is that the closed forms the symbolic engine returns for the
three cumulants and the adjacent-face cube completion \emph{sum to} the
corpus's coefficient --- a polynomial identity between the engine's
output and the corpus's formula. It is not a Lean derivation of the
cumulants, and the Lean source says so in its own docstring. This
document repeats the qualification because an internal summary of the
same fact dropped it.

Every eigencomponent of every resolvent is verified to be an exact
eigenvector of every single-link $H_0$, symbolically --- a check that
has no per-rank analogue, because at a fixed rank a coincidence of
numbers is indistinguishable from an identity.

\subsection{Why this is the right shape of argument}

The methodological content generalizes beyond this program.

An identity between rational functions can be established in two ways:
verify it at enough points and bound the degrees, or compute in the
function field. The first requires an argument about the size of the
objects, which is a second piece of mathematics and is where the
difficulty hides. The second requires only that the algorithm's
dependence on the parameter be rational --- which is a property one can
\emph{inspect}, predicate by predicate, rather than estimate.

The title of the decision record is the summary: the degree bound is
proved by removing it.

\begin{scopenote}
The symbolic engine is valid at ranks where the generic partition
labelling applies. Exceptional low ranks, where representations
degenerate, are handled separately and are the reason $\SU(2)$ and in
places $\SU(4)$ carry their own entries in the register.

This degree bound must not be confused with the \emph{other} degree
bound in this program's history --- the proposed kinematic mechanism
for the fourth-order tier collapse, which was retracted
(\S\ref{sec:fourth}, \S\ref{sec:register}). That one was withdrawn
because it was wrong. This one was removed because it became
unnecessary.
\end{scopenote}
% ---- end 15_symbolic.tex -------------------------------------------
% ---- begin 15_sixth_order.tex --------------------------------------
\section{Sixth order: a shape that cannot be removed}
\label{sec:sixth}

The fourth-order shape data closes in a five-element span. The question
this section answers is whether the sixth order stays inside it. It does
not, and the proof is a divisibility argument that needs no knowledge of
the sixth-order dynamics at all.

\subsection{Conventions}

In the Bloch variables $z_j$ write
\begin{equation}
  a_j = 2 - z_j - z_j^{-1},
  \qquad q = \sum_j a_j,
  \qquad e_2 = \sum_{i<j} a_i a_j,
  \qquad e_3 = a_1 a_2 a_3 ,
\end{equation}
and let $U = \psi\psi^{*}$ with $\psi^{*}\psi = q$, $P_c = U/q$,
$Q_c = I - P_c$, $S = L_{\downarrow} - 4I$, and $R$ the cross-plane part
of $S$.

Two warnings, both of which have caused errors here. The invariants
$e_2, e_3$ are Bloch invariants, \emph{not} representation Casimirs;
and $U$ is the up-Laplacian, not its normalized projector, so its kernel
is the complementary plane rather than the carrier line. Cleared
identities hold at $q = 0$; normalized carrier formulas require
$q > 0$.

\subsection{The complete folded operator}

Let $P_E$ be the isolated electric model-space projector,
$D = Q_E/(E_0 - H_0)$, and $V_0$ the perturbation. If
$P_E V_0 P_E = a P_E$, set $V = V_0 - aI$ \emph{on the whole Hilbert
space}: shifting the complementary block as well is essential. With
$\chi_0 = P_E$, intermediate normalization gives the exact recursions
\begin{equation}
  K_n = P_E V \chi_{n-1},
  \qquad
  \chi_n = D V \chi_{n-1} - D \sum_{j=1}^{n-1} \chi_{n-j} K_j ,
\end{equation}
and with $M = P_E + \chi^{*}\chi$ the canonical Hermitian effective
operator is $M^{1/2} K M^{-1/2}$. Writing
\begin{equation}
  A_m = P_E V (DV)^{m-1} P_E,
  \qquad
  C_2 = P_E V D^3 V P_E,
\end{equation}
and $B_m$ for the sum of $A_m$ with exactly one of its $m-1$ slots
replaced by $D^2$:

\begin{theorem}[The eighteen-word sixth-order fold]
\label{thm:folds}
\begin{equation}
  H_6 = A_6 - \tfrac12\{A_4,B_2\} - \tfrac12\{A_3,B_3\}
        - \tfrac12\{A_2,B_4\} + \tfrac12\{A_2^2,C_2\}
        + \tfrac38\{A_2,B_2^2\} + \tfrac14 B_2 A_2 B_2 .
  \label{eq:F6}
\end{equation}
Expanded, this is exactly eighteen words in $P_E, V, D$. The same
recursion reproduces the established $H_4 = A_4 - \tfrac12\{A_2,B_2\}$.
\end{theorem}
\gid{RESULT:G9\_HERMITIAN\_FOLDED\_SUPPORT}

The $\{A_3,B_3\}$ term is present and it matters: dropping odd electric
histories gives an incomplete formula. No odd order is dropped anywhere
in the recursion. Independent rational matrix checks verify both the
invariant-subspace equation and the normalization through degree six,
with a rank-two $P_E$ and non-commuting effective coefficients.

\notasserted{Equation~\eqref{eq:F6} is the universal folded operator
support in this convention. It does not evaluate Wilson plaquette matrix
elements, establish a specific support, or carry out connected
multi-plaquette rooted subtraction. Those require the actual electric
and Haar implementation on each support.}

\subsection{Exact word symbols}

In a polynomial carrier gauge, exact symmetrization gives each word its
cleared symbol in $\Q[q, e_2, e_3]$. The table below is computed, not
inferred, and all forty words of length zero through three agree
coefficient by coefficient with independently implemented spatial
Laurent operators.

\begin{center}
\begin{tabular}{@{}ll@{\qquad}ll@{}}
\toprule
Word & $\sigma(W)$ & Word & $\sigma(W)$ \\
\midrule
$I$        & $q$          & $RR$        & $q e_2 + 3 e_3$ \\
$U$        & $q^2$        & $RUR$       & $4 e_2^2$ \\
$S$        & $-4q$        & $RSR$       & $(q-4)(q e_2 + 3e_3) - 4e_2^2$ \\
$S^2$      & $16q$        & $RRR$       & $-2e_2^2 - 2q e_3$ \\
$R$        & $-2e_2$      & $URR,\,RRU$ & $q^2 e_2 + 3 q e_3$ \\
$UR,\,RU$  & $-2q e_2$    &             & \\
\bottomrule
\end{tabular}
\end{center}

A word census must retain coefficients and combine them before testing
shape membership: $\sigma(RUR + 2\,RRR) = -4q e_3$, although both
individual words contain $e_2^2$. Support is not unique modulo operator
identities, and the invariant object being tested is the complete
carrier polynomial. Missing direct input is represented as unknown, not
as zero.

\subsection{The band coefficient at sixth order}

At fixed nonzero momentum, with scalar anchors removed,
\begin{equation}
  H_{\mathrm{eff}}(u) = u^2 t_3\, q\, Q_c + u^3 H_3 + u^4 H_4
    + u^5 H_5 + u^6 H_6 + O(u^7).
\end{equation}
The recorded $\SU(3)$ third-order factorization gives
$Q_c H_3 P_c = 0$, so the off-diagonal block first appears at order
$u^2$ in the scaled operator, from $H_4$. Expanding the Schur complement
for the scalar carrier branch,
\begin{equation}
  \lambda_6 = \frac{\sigma(H_6)}{q}
    - \frac{q\,\sigma(H_4^2) - \sigma(H_4)^2}{t_3\, q^3}.
  \label{eq:S6}
\end{equation}
This is not an assumption that $H_5$ vanishes. $H_5$ first couples to
$H_4$ at order $u^7$ in $H_{\mathrm{eff}}$, as does the $uH_3$
correction to the complementary inverse, and unknown off-diagonal $H_6$
contributes later still. The decoupling of $H_3$ is an explicit recorded
input rather than an assumption, and any scalar vacuum subtraction is
part of $H_6$ and changes nothing.

Evaluating all twenty-five ordered pairs from the actual $H_4$ support
$\{I,U,S,S^2,R\}$, including their projected subtraction, gives zero for
twenty-four of them. Only $RR$ survives, in agreement with direct
composition of the full assembled spatial kernel. With
\begin{equation}
  \kappa = \frac{4C_{\mathrm{shp}}^2}{t_3}
  = \frac{169924002729806205028788409}{619343593697933825385600000} > 0,
  \qquad t_3 = \tfrac{5}{612},
\end{equation}
the known dynamical mixing support is
$-(\kappa/q)\,RR + (\kappa/q^2)\,RUR$, with normalized expectation
\begin{equation}
  \lambda_6^{\mathrm{mix}}
  = -\kappa\left(\frac{e_2}{q} + \frac{3e_3}{q^2}
      - \frac{4e_2^2}{q^3}\right).
\end{equation}

\notasserted{The $RUR$ here arises from inserting $P_c = U/q$ in the
projected subtraction. It is not evidence that a direct six-insertion
electric history carries the local word $RUR$, and the ratio $1:3:-4$
belongs to this mixing term rather than to the full coefficient.}

\subsection{Non-cancellation}

\begin{theorem}[An extra rational shape survives every local direct term]
\label{thm:noncancellation}
Assume the order-six electric-shell effective operator is a
finite-range, translation-covariant stencil with Laurent-polynomial
entries in this basis, and write $F = \sigma(H_6)$. Then \eqref{eq:S6}
gives
\begin{equation}
  q^3 \lambda_6 = q^2 F - \kappa\big(q^2 e_2 + 3q e_3 - 4e_2^2\big),
  \label{eq:N6}
\end{equation}
and $q$ does not divide the right-hand side, for \emph{any} Laurent
polynomial $F$.
\end{theorem}

\begin{proof}
Modulo $q$, the right-hand side of \eqref{eq:N6} is $4\kappa e_2^2$.
Evaluate at
\begin{equation}
  z_1 = z_2 = -1, \qquad z_3 = 5 + 2\sqrt{6},
\end{equation}
all nonzero, so the point lies in the Laurent ring. Since
$z_3^{-1} = 5 - 2\sqrt6$, we get $a = (4, 4, -8)$, hence $q = 0$ and
$e_2 = -48$. The value there is $4\kappa \cdot 48^2 = 9216\,\kappa \neq 0$,
and a Laurent polynomial divisible by $q$ would vanish at that point.
\end{proof}
\gid{RESULT:G9\_LOCAL\_SIXTH\_NONCANCELLATION}

\noindent\Tone{}. Consequently $\lambda_6$ cannot lie in the
fourth-order shape span $\{1, q, e_2, 4e_2/q, e_3/q\}$, and no local
direct $RUR$ term removes the extra rational shape: a direct $h\,RUR$
contributes $4h e_2^2/q$, which cannot cancel an $e_2^2/q^3$ component.
For symmetric $F$ the remainder modulo $q^2$ is exactly
$4\kappa e_2^2 - 3\kappa q e_3$, so both $e_3/q^2$ and $e_2^2/q^3$ are
present. The $e_2/q$ term alone \emph{can} be cancelled, by
$F = \kappa e_2$ --- for instance a local $H_6 = -(\kappa/2)R$ --- so
the $1:3:-4$ ratio is not a constraint on all sixth-order shapes.

Note what the proof does not use: any knowledge of $F$. The extra shape
is forced by the fourth-order data alone, through the Schur complement,
whatever the sixth-order dynamics turns out to be.

\subsection{A withdrawn derivation}

An earlier version of this work enumerated closed rooted walks --- $90$
four-step unit-edge walks, of which $66$ retrace immediately and $24$
are planar squares, together with $192$ non-retracing six-step walks
using all three axes --- and read a sixth-order amplitude off that
census.

The derivation is withdrawn. Edge walks are not ordered plaquette
insertions with electric states, Haar weights, projectors and folds, and
no map between them was supplied; the census therefore determines no net
coefficient of any of the candidate words. The old amplitude formula was
assigned rather than derived. The old carrier function used
$\sigma(R) = 0$ and $\sigma(U) = q$, where the actual unnormalized
symbols are $-2e_2$ and $q^2$; it returned a nonzero expression for a
defect it separately reported as zero. The implementation now computes
the symbols and tests the defect, and the old interface raises an
explicit missing-dynamics error rather than returning a number.

What is withdrawn is a claimed derivation, not the possibility of a
physical amplitude of that shape. Theorem~\ref{thm:noncancellation} is
independent of the census entirely. Computing the direct term on a
connected multi-plane support is Problem~\ref{prob:g9}, and it is the
cheapest decisive computation currently open in this program.
% ---- end 15_sixth_order.tex ----------------------------------------
% ---- begin 16_wick_shell.tex ---------------------------------------
\section{A rank-uniform local shell construction}
\label{sec:wick}

This section is a separate result in its own regime, included because it
is exact, rank-uniform, and independent of everything above --- not
because it continues it. The object is the formal local Wilson
oscillator expansion around a single plaquette, and the point is that
its spectrum admits a finite exact recursion at every perturbative
order with coefficients in $\Q[N, N^{-1}]$, needing no Gram-matrix
inversion.

Conflating this construction with the flux band of
\S\ref{sec:carrier}--\S\ref{sec:sixth}, or with the Wilson transfer
operator of \S\ref{sec:wilson}, would be an error. They are four
different objects --- the physical $\SU(N)$ transfer band, the $\SU(3)$
internal Hamiltonian carrier, the conditional block, and this local
oscillator --- and no cross-regime continuation between them is
asserted anywhere in this document.

\subsection{The exact algebra}

Let $X$ be a traceless Hermitian $N \times N$ matrix with density
proportional to $e^{-\Tr(X^2)}$, write $P_k = \Tr(X^k)$ with $P_0 = N$
and $P_1 = 0$, and use a traceless Hermitian basis normalized by
$\Tr(T_a T_b) = \delta_{ab}$. The covariance is
\begin{equation}
  \mathbb{E}\big[X_{ab} X_{cd}\big]
  = \tfrac12\left(\delta_{ad}\delta_{bc}
    - \tfrac{1}{N}\delta_{ab}\delta_{cd}\right).
\end{equation}
With $D = \tfrac12\Delta$ and $E$ the polynomial Euler operator, the
completeness identity for this basis gives the exact relations
\begin{align}
  D P_k &= \frac{k}{2}\sum_{l=0}^{k-2} P_l P_{k-2-l}
           - \frac{k(k-1)}{2N} P_{k-2},
  \\
  D(fg) &= (Df)g + f(Dg) + \nabla f \cdot \nabla g,
  \\
  \nabla P_k \cdot \nabla P_l
        &= k\,l\left(P_{k+l-2} - \frac{1}{N} P_{k-1}P_{l-1}\right).
\end{align}
These are exact at every integer rank, including ranks at which some
traces become dependent --- which is the property that makes the
construction rank-uniform rather than rank-by-rank. For homogeneous $F$
of positive degree $d$, Gaussian integration by parts gives
$\mathbb{E}_\mu[F] = \mathbb{E}_\mu[DF]/d$, and with
$\mathbb{E}_\mu[1] = 1$ and vanishing odd moments this is a finite exact
moment recursion.

\subsection{Shell projectors without a Gram inverse}

\begin{proposition}[Terminating Wick conjugation]
\label{prop:wick}
Put $L = E - D$. Since $[E, D] = -2D$, the terminating exponential
$W = \exp(-D/2)$ satisfies
\begin{equation}
  L\,W = W\,E .
\end{equation}
Hence, with $H_d$ the extraction of the homogeneous degree-$d$ part,
\begin{equation}
  \Pi_d = \exp(-D/2)\, H_d \,\exp(D/2)
\end{equation}
is the exact oscillator-shell projector, and on each finite polynomial
space $R_s = \sum_{d \neq s} \Pi_d/(s-d)$ is the exact shell resolvent.
\end{proposition}

The content of Proposition~\ref{prop:wick} is that $L$ and $E$ are
conjugate by a \emph{terminating} exponential, so the shell
decomposition of the interacting operator is the shell decomposition of
the Euler operator, transported. No Gram matrix is formed and none is
inverted, which is precisely the step that makes rank-uniformity
delicate elsewhere in this program.

\subsection{The coefficients}

\begin{theorem}[All-rank local coefficients]
\label{thm:wickcoeffs}
The terminating conjugation gives exact shell projectors and a unique
all-order simple-branch formal recursion in $\Q[N, N^{-1}]$, with state
degree at most $s + 4r$. The two formal local gaps have exact all-rank
coefficients through $\beta^{-2}$, including $c_3^{-}$ and
$c_4^{\pm}$. The five corrections $c_0, \ldots, c_4$ are strictly
negative at every allowed rank, by positive shifted numerator
polynomials.
\end{theorem}
\gid{RESULT:LOCAL\_CLASS\_WICK\_RECURRENCE, RESULT:LOCAL\_CLASS\_WICK\_COEFFICIENTS, RESULT:LOCAL\_CLASS\_WICK\_SIGNS}

The construction applies to the vacuum, the first charge-even shell for
$N \ge 2$, and the first charge-odd shell for $N \ge 3$. It recovers the
previously archived $c_0$, $c_1$, $c_2$ formulas in both sectors and the
archived all-rank $c_3^{+}$, and additionally derives $c_3^{-}$ and both
all-rank $c_4$ formulas. The certificate carries full polynomial
eigen-equation residuals rather than only agreement with the archived
constants --- a distinction worth stating, because agreement with a
previous value is a weaker check than satisfying the equation.

The recovery is by an independent route: the archive engine uses
full-GUE cut-and-join and then subtracts the trace, while this one
applies $D$ directly in traceless coordinates.

\subsection{A companion algebraic identity}

\begin{theorem}[Shifted-power carrier symbols]
\label{thm:rsmr}
For the cubic Laurent operators with $S = L_{\downarrow} - 4I$,
\begin{equation}
  \sigma(R S^m R) = (q-4)^m\big(q e_2 + 3 e_3\big) - 4\,\Pi_m(q)\, e_2^2
\end{equation}
for every natural $m$, with $\Pi$ defined polynomially also at $q = 0$.
\end{theorem}
\gid{RESULT:RSM\_R\_CARRIER\_SYMBOL\_MASTER\_THEOREM}

\noindent This closes the word table of \S\ref{sec:sixth} at all powers,
and is formalized in Lean as an all-power operator statement.

\begin{scopenote}
Theorem~\ref{thm:wickcoeffs} concerns the \emph{formal} local class
oscillator only. No error bound for the compact-group eigenvalues is
claimed, and local spectral coefficients do not by themselves identify
the interacting Wilson spectrum, an interacting-volume estimate, or a
continuum source measure. Theorem~\ref{thm:rsmr} is an algebra identity
and carries no sixth-order dynamical population claim.
\end{scopenote}
% ---- end 16_wick_shell.tex -----------------------------------------
% ---- begin 17_external.tex -----------------------------------------
\section{External comparison}
\label{sec:external}

Published work enters this program as evidence rather than as authority.
A paper is unchecked until something checks it, exactly like an internal
document; the reason an external agreement counts for more is not that
it is published but that it was produced without knowledge of this
program, so agreement is not bookkeeping.

The map holds \LitPapers{} indexed papers and \LitEdges{} recorded
relations to named claims here. \LitVerified{} of those edges are
verified --- the paper was obtained, read, and compared against the
claim --- across \LitObtained{} distinct papers. The remainder are
recorded targets for papers not yet obtained, which is itself useful: a
paper heavily cited within this web and still unobtained is
automatically the next acquisition target.

An entry earns its place by naming a specific claim and saying what
relationship it has to it. \texttt{workhouse lit --for G19} answers the
query in the other direction, and validation rejects an entry whose
target does not resolve.

\subsection{A 1989 table, and rationals computed cold}

The sharpest external check available to this program is a
Hamiltonian strong-coupling series for $\SU(3)$ glueball masses
published in 1989~\cite{HAMER_1989}, whose Table~1 is pinned here by
digest. Under the normalization bridge $m_n = 2^{\,n-1} a_n$, its
coefficients are compared against the $\Gamma$-point coefficients this
program derives as exact rationals.

In the $C$-odd ($1^{+-}$) channel --- the sector of the carrier, whose
series is \eqref{eq:thirdorder} --- the agreement runs through third
order:

\begin{center}
\begin{tabular}{@{}l l l l@{}}
\toprule
Order & Published & This program & Gap (bound) \\
\midrule
$n=0$ & $\tfrac{16}{3}/2$ & $\tfrac83 = E_{\mathrm{flat}}(0)$ & exact \\
$n=1$ & $a_1 = +1$ & the $u$ coefficient $+1$ & exact \\
$n=2$ & $2a_2 = 0.0359477124184$ & $\tfrac{11}{306}$
      & $9.94\times10^{-14}$ ($2\times10^{-13}$) \\
$n=3$ & $4a_3 = -0.437135556836$ & $-\tfrac{109151}{249696}$
      & $1.11\times10^{-12}$ ($3\times10^{-12}$) \\
\bottomrule
\end{tabular}
\end{center}

\noindent In the $C$-even ($0^{++}$) channel, $2a_2 = -0.21274509804$
against $-\tfrac{217}{1020}$ with gap $7.84\times10^{-13}$, and
$4a_3 = -0.1039004229144$ against $-\tfrac{54049}{520200}$ with gap
$1.36\times10^{-13}$. The sign structure matches as well: $a_1 = -1$ in
the scalar channel, where the carrier's is $+1$.

Each bound above is the printed half-ulp of the published table scaled
with margin, so these are agreements to the precision at which the 1989
table was printed --- not approximate corroborations. The rationals
$11/306$ and $109151/249696$ are the third-order coefficients of
$E_{\mathrm{flat}}$ in \eqref{eq:thirdorder}, computed here from Haar
integration and the cellular complex with no knowledge of the table.

At fourth order, $8a_4 = -0.7751458630184$ against
$m^{(4)}_{\Gamma} = -0.7751458630189173$, gap $5.17\times10^{-13}$
against a bound of $5.3\times10^{-13}$.

\begin{scopenote}
These are \Ttwo{} checks: float agreement within a stated tolerance, to
the precision of a printed table. They corroborate the coefficients;
they do not prove them, and the exact statements of
\S\ref{sec:second}--\S\ref{sec:grading} do not rest on them. Note also
that the pinned revision of the source compares only $2a_2$ and $4a_3$
in its own text --- the fourth-order agreement is registered here and
not leaned on there.
\end{scopenote}

\subsection{The two contradicting edges}

Two edges in the map are recorded as \emph{contradicts}, both from a
1976 strong-coupling series, and both against this program: one against
the third-order $C$-odd coefficient, one against the $\Gamma$-point
fourth-order value. They originate in a disagreement the 1989 paper
itself reports at orders $x^3$ and $x^4$ with the series it supersedes.

Both are recorded as not yet obtained. Until that primary source is
read, this program has two unresolved external contradictions against
coefficients it otherwise regards as settled, and says so rather than
resolving them by the age of the sources. Obtaining that paper is
listed with the acquisition targets below.

\subsection{How far a novelty claim here can be trusted}

Not as far as the size of the index suggests, and the limitation is
worth stating precisely rather than as a disclaimer.

Only \LitVerified{} of \LitEdges{} recorded edges are verified. The
unverified remainder is not a random sample: it includes several of the
\emph{nearest} neighbours to this program's own results. The main
Nuclear Physics B paper of the 1981 strong-coupling glueball series is
recorded as not obtained --- only its errata were --- and the 1974 and
1979 large-$N$ papers that bear most directly on the rank grading of
\S\ref{sec:grading} are likewise unobtained, with the ledger noting for
one of them that its definitions were used from memory and only their
sign consequences checked.

Consequently: where this document says a result is novel, that means
\emph{no verified entry in this index contains it}, which is a weaker
statement than priority. Two of the strongest results are affected. The
rank-grading cancellation of \S\ref{sec:grading} sits closest to
large-$N$ counting arguments in papers this program has not read; and
the fixed-spacing transfer construction of \S\ref{sec:wilson} sits
alongside the classical positivity and self-adjoint transfer-matrix
results of Osterwalder and Seiler, which the index does register.

Acquiring the unobtained near neighbours is cheap relative to
everything else in \S\ref{sec:obligations}, and it is the one item on
that list that could \emph{subtract} from the program rather than add
to it. That is a reason to do it, not a reason to defer it.

% ---- begin gen_literature.tex --------------------------------------
% Generated by paper/master/generate_apparatus.py.  Do not edit.
\begin{longtable}{@{}l l l l l@{}}
\caption{Verified external-evidence edges. Each row is a published result that was obtained, read, and checked against a named claim of this program. Relation is the ledger's own classification; \emph{contradicts} and \emph{confusable} are listed first because they are the rows that cost the program something.}\label{tab:literature}\\
\toprule
Paper & Ref. & Year & Bears on & Relation \\
\midrule
\endfirsthead
\toprule Paper & Ref. & Year & Bears on & Relation \\ \midrule \endhead
\bottomrule
\endfoot
\texttt{\small HSB\_\allowbreak{}2000} & \cite{HSB_2000} & 2000 & \texttt{HAMER\_\allowbreak{}A4\_\allowbreak{}NUM} & confusable with \\
\texttt{\small HAMER\_\allowbreak{}1989} & \cite{HAMER_1989} & 1989 & \texttt{BAND\_\allowbreak{}EVEN\_\allowbreak{}BOTTOM} & corroborates \\
\texttt{\small HAMER\_\allowbreak{}1989} & \cite{HAMER_1989} & 1989 & \texttt{C1} & corroborates \\
\texttt{\small LLL\_\allowbreak{}2006} & \cite{LLL_2006} & 2006 & \texttt{C1} & corroborates \\
\texttt{\small HAMER\_\allowbreak{}1989} & \cite{HAMER_1989} & 1989 & \texttt{D\_\allowbreak{}3} & corroborates \\
\texttt{\small SCHIERHOLZ\_\allowbreak{}1988} & \cite{SCHIERHOLZ_1988} & 1988 & \texttt{G18} & corroborates \\
\texttt{\small HAMER\_\allowbreak{}1989} & \cite{HAMER_1989} & 1989 & \texttt{M3\_\allowbreak{}EVEN\_\allowbreak{}K0} & corroborates \\
\texttt{\small KPS\_\allowbreak{}1981} & \cite{KPS_1981} & 1981 & \texttt{R14} & corroborates \\
\texttt{\small CBB\_\allowbreak{}2026} & \cite{CBB_2026} & 2026 & \texttt{R2} & corroborates \\
\texttt{\small KPS\_\allowbreak{}1981} & \cite{KPS_1981} & 1981 & \texttt{G7} & supplies a value \\
\texttt{\small HAMER\_\allowbreak{}1989} & \cite{HAMER_1989} & 1989 & \texttt{HAMER\_\allowbreak{}A4\_\allowbreak{}NUM} & supplies a value \\
\texttt{\small MUNSTER\_\allowbreak{}1985\_\allowbreak{}TM} & \cite{MUNSTER_1985_TM} & 1985 & \texttt{C2} & supplies a comparison \\
\texttt{\small CAO\_\allowbreak{}2023} & \cite{CAO_2023} & 2023 & \texttt{C7} & supplies a comparison \\
\texttt{\small LEMOINE\_\allowbreak{}2026} & \cite{LEMOINE_2026} & 2026 & \texttt{C7} & supplies a comparison \\
\texttt{\small HAZRA\_\allowbreak{}2026} & \cite{HAZRA_2026} & 2026 & \texttt{G14} & supplies a comparison \\
\texttt{\small AT\_\allowbreak{}2020\_\allowbreak{}SU3} & \cite{AT_2020_SU3} & 2020 & \texttt{G18} & supplies a comparison \\
\texttt{\small BESIII\_\allowbreak{}2026} & \cite{BESIII_2026} & 2026 & \texttt{G18} & supplies a comparison \\
\texttt{\small CHEN\_\allowbreak{}2006} & \cite{CHEN_2006} & 2006 & \texttt{G18} & supplies a comparison \\
\texttt{\small LLL\_\allowbreak{}2006} & \cite{LLL_2006} & 2006 & \texttt{G18} & supplies a comparison \\
\texttt{\small MORNINGSTAR\_\allowbreak{}2025} & \cite{MORNINGSTAR_2025} & 2025 & \texttt{G18} & supplies a comparison \\
\texttt{\small MP\_\allowbreak{}1999} & \cite{MP_1999} & 1999 & \texttt{G18} & supplies a comparison \\
\texttt{\small AT\_\allowbreak{}2021\_\allowbreak{}SUN} & \cite{AT_2021_SUN} & 2021 & \texttt{G19} & supplies a comparison \\
\texttt{\small BB\_\allowbreak{}1983} & \cite{BB_1983} & 1983 & \texttt{G19} & supplies a comparison \\
\texttt{\small FLPS\_\allowbreak{}2026} & \cite{FLPS_2026} & 2026 & \texttt{G19} & supplies a comparison \\
\texttt{\small URZUA\_\allowbreak{}2019} & \cite{URZUA_2019} & 2019 & \texttt{G19} & supplies a comparison \\
\texttt{\small BORGA\_\allowbreak{}2024} & \cite{BORGA_2024} & 2024 & \texttt{G5} & supplies a comparison \\
\texttt{\small CAO\_\allowbreak{}2023} & \cite{CAO_2023} & 2023 & \texttt{G5} & supplies a comparison \\
\texttt{\small OBZ\_\allowbreak{}1985} & \cite{OBZ_1985} & 1985 & \texttt{G5} & supplies a comparison \\
\texttt{\small AT\_\allowbreak{}2021\_\allowbreak{}SUN} & \cite{AT_2021_SUN} & 2021 & \texttt{G6} & supplies a comparison \\
\texttt{\small CM\_\allowbreak{}2003} & \cite{CM_2003} & 2003 & \texttt{G6} & supplies a comparison \\
\texttt{\small BALAJI\_\allowbreak{}2026} & \cite{BALAJI_2026} & 2026 & \texttt{R2} & supplies a comparison \\
\texttt{\small CKS\_\allowbreak{}2021} & \cite{CKS_2021} & 2021 & \texttt{U1} & supplies a comparison \\
\texttt{\small DRS\_\allowbreak{}2021} & \cite{DRS_2021} & 2021 & \texttt{U1} & supplies a comparison \\
\texttt{\small KRS\_\allowbreak{}2023} & \cite{KRS_2023} & 2023 & \texttt{U1} & supplies a comparison \\
\texttt{\small SZH\_\allowbreak{}1997} & \cite{SZH_1997} & 1997 & \texttt{U1} & supplies a comparison \\
\texttt{\small CS\_\allowbreak{}2006} & \cite{CS_2006} & 2006 & \texttt{C7} & supplies a method \\
\texttt{\small JW\_\allowbreak{}2006} & \cite{JW_2006} & 2006 & \texttt{G17} & supplies a method \\
\texttt{\small UELTSCHI\_\allowbreak{}2004\_\allowbreak{}CLUSTER} & \cite{UELTSCHI_2004_CLUSTER} & 2004 & \texttt{G17} & supplies a method \\
\texttt{\small YAROTSKY\_\allowbreak{}2004\_\allowbreak{}GS} & \cite{YAROTSKY_2004_GS} & 2004 & \texttt{G17} & supplies a method \\
\texttt{\small BB\_\allowbreak{}1983} & \cite{BB_1983} & 1983 & \texttt{G18} & supplies a method \\
\texttt{\small JW\_\allowbreak{}2006} & \cite{JW_2006} & 2006 & \texttt{G18} & supplies a method \\
\texttt{\small NSY\_\allowbreak{}2019} & \cite{NSY_2019} & 2019 & \texttt{G18} & supplies a method \\
\texttt{\small YAROTSKY\_\allowbreak{}2004\_\allowbreak{}QP} & \cite{YAROTSKY_2004_QP} & 2004 & \texttt{G18} & supplies a method \\
\texttt{\small JW\_\allowbreak{}2006} & \cite{JW_2006} & 2006 & \texttt{G19} & supplies a method \\
\texttt{\small KLEIN\_\allowbreak{}ROSENBERGER\_\allowbreak{}2020\_\allowbreak{}WKB} & \cite{KLEIN_ROSENBERGER_2020_WKB} & 2020 & \texttt{G19} & supplies a method \\
\texttt{\small SCHIERHOLZ\_\allowbreak{}1988} & \cite{SCHIERHOLZ_1988} & 1988 & \texttt{G19} & supplies a method \\
\texttt{\small UELTSCHI\_\allowbreak{}2004\_\allowbreak{}CLUSTER} & \cite{UELTSCHI_2004_CLUSTER} & 2004 & \texttt{G19} & supplies a method \\
\texttt{\small YAROTSKY\_\allowbreak{}2004\_\allowbreak{}GS} & \cite{YAROTSKY_2004_GS} & 2004 & \texttt{G19} & supplies a method \\
\texttt{\small YAROTSKY\_\allowbreak{}2004\_\allowbreak{}QP} & \cite{YAROTSKY_2004_QP} & 2004 & \texttt{G19} & supplies a method \\
\texttt{\small JW\_\allowbreak{}2006} & \cite{JW_2006} & 2006 & \texttt{G20} & supplies a method \\
\texttt{\small JW\_\allowbreak{}2006} & \cite{JW_2006} & 2006 & \texttt{G22} & supplies a method \\
\texttt{\small JW\_\allowbreak{}2006} & \cite{JW_2006} & 2006 & \texttt{G23} & supplies a method \\
\texttt{\small MUNSTER\_\allowbreak{}1985\_\allowbreak{}TM} & \cite{MUNSTER_1985_TM} & 1985 & \texttt{G3} & supplies a method \\
\texttt{\small KRS\_\allowbreak{}2023} & \cite{KRS_2023} & 2023 & \texttt{R2} & supplies a method \\
\end{longtable}

\nocite{*}

% ---- end gen_literature.tex ----------------------------------------
% ---- end 17_external.tex -------------------------------------------

\appendix
% ---- begin A0_reproducibility.tex ----------------------------------
\section{Reproducibility}
\label{app:repro}

Every claim in this document that carries a \Tzero{}, \Tone{} or
\Ttwo{} badge can be re-established from a clean checkout in seconds to
minutes. This appendix says how, and --- more usefully --- says what each
kind of re-establishment does and does not prove.

\subsection{From a clean checkout}

\begin{verbatim}
git clone https://github.com/ats314/WORKHOUSE
cd WORKHOUSE
uv sync --all-extras --frozen
uv run --no-sync workhouse --help
\end{verbatim}

Then:

\begin{verbatim}
make verify        # every registered check: T1 and T2, a few seconds
make lean          # T0: proof-check the Lean core (requires elan)
make status        # the contradiction and gap registers
workhouse verify --tier 1              # only the exact re-derivations
workhouse verify -v                    # every tolerance, printed
workhouse verify --only 'h_4^side'     # one claim, with its numbers
\end{verbatim}

A single claim is the useful unit. \texttt{workhouse verify -{}-only} takes a
claim name and prints the numbers it establishes, which is what makes a
disagreement arguable rather than merely reported. A certification a
reader cannot cheaply reproduce is a claim of authority, not evidence,
and the design of this repository takes that seriously enough to attach
the reproducing command to every row of \srcpath{CERTIFIED.md}.

\subsection{Querying a claim}

\begin{verbatim}
workhouse why C2                   # everything recorded about one id
workhouse why t_N                  # a constant: value, formalizations, sources
workhouse branches C2              # both sides of a dispute, side by side
workhouse derive C2 G3 --out f.md  # evidence chain, registered edges only
workhouse search '5/612'           # search by exact VALUE, not by name
workhouse lit --for G19            # published work bearing on a claim
workhouse brief G19 --json         # target snapshot with freshness
\end{verbatim}

\texttt{search} matches by value rather than spelling, which matters
because the join keys of this corpus are exact rationals and no
natural-language query retrieves $109151/249696$. It also carries two
warnings a text search cannot: names the register forbids because they
regenerate a dissolved contradiction, and names coined in this program
that its own source corpus never uses.

\subsection{What each tier actually establishes}

\begin{description}[leftmargin=1.2em,style=nextline]
\item[\Tzero{} --- Lean 4, no \texttt{sorry}.]
  The statement is checked by the Lean kernel from explicit definitions,
  with axioms restricted to \texttt{propext}, \texttt{Classical.choice}
  and \texttt{Quot.sound}. The axiom set is exported and validated rather
  than asserted: \srcpath{scripts/export\_lean\_dependencies.py} builds
  strictly and records elaborated constants and transitive axioms into
  \srcpath{ledger/lean\_dependencies.json}, and the loader checks these
  against declaration records and source fingerprints. Identifiers
  mentioned in comments are not proof dependencies.

  What a T0 badge does \emph{not} establish is importance. Several
  registered theorems formalize arithmetic evaluations of named
  rationals. The registry records for each theorem whether it proves an
  entire named check (\texttt{promotes}) or one ingredient of a larger
  statement (\texttt{lean\_support}); supporting a statement never
  silently promotes the whole statement's tier.

\item[\Tone{} --- exact re-derivation.]
  The statement is recomputed symbolically in exact arithmetic from
  stated definitions. Corpus rationals are \texttt{sympy.Rational};
  values the corpus records only as floats carry a \texttt{\_NUM} suffix
  and are Python floats. A float that reads as exact is the most
  dangerous defect available in this repository, and the boundary is
  guarded by a dedicated test.

\item[\Ttwo{} --- numerical agreement.]
  Float agreement within a tolerance the check itself prints. A T2 pass
  is that and nothing more. Tolerances are never widened to make a
  disagreement disappear; a disagreement is registered as an explicit
  finding that asserts the discrepancy.

\item[\Tthree{} --- no whole-statement machine certification.]
  The default. An analytic result may be fully proved and still be T3:
  the badge reports machine coverage, not mathematical status. Status,
  evidence and tier are three independent fields in
  \srcpath{ledger/results.yaml} for exactly this reason.
\end{description}

\subsection{The one inference this apparatus does not license}

A finite check certifies the statement it checks. It does not certify a
general theorem of which that statement is an instance, and the converse
also holds: a complete analytic proof is not weakened by absent Lean
coverage.

This cuts both ways in practice, and both directions have caused errors
here. A suite of exact checks passing at ranks $N = 3, 4, 5, 6$ is not a
statement about all $N$ --- which is why the symbolic-$N$ engine of
\S\ref{sec:symbolic} exists, and why its degree bound was retracted when
one projection turned out to be uncounted. Conversely, a separate Lean
library in this program's archive comprises seventy-one modules with no
\texttt{sorry} and was found on audit to prove mostly tautologies; a
theorem count is not a measure of coverage, and no count in this
document should be read as one.

\subsection{Regenerating this document}

\begin{verbatim}
python paper/master/generate_apparatus.py
tectonic paper/master/workhouse_master.tex        # the monograph
tectonic paper/master/workhouse_band_paper.tex    # Part I, standalone
\end{verbatim}

The generator reads \srcpath{literature/index.yaml} and
\srcpath{ledger/\{gaps,results,theorems,contradictions\}.yaml} and writes
the bibliography, the external-evidence table, the route register, the
results table and the count macros. Every number the prose quotes is one
of those macros. If a count here disagrees with the repository, the
generator has not been re-run; the repository is right.
% ---- end A0_reproducibility.tex ------------------------------------
% ---- begin A1_verification.tex -------------------------------------
\section{Verification index}
\label{app:verification}

Table~\ref{tab:results} lists every registered analytic result with its
recorded status and evidence level. The two columns are independent
fields and are meant to be read against each other: \texttt{proven} with
\texttt{record-backed} evidence means the argument exists and the machine
artefact does not.

Of \ResTotal{} registered results, \ResProven{} are recorded
\texttt{proven}, \ResConditional{} \texttt{conditional},
\ResDisputed{} \texttt{disputed} and \ResFalsified{} \texttt{falsified}.
The Lean tree carries \LeanTotal{} registered theorems with
\LeanSorry{} occurrences of \texttt{sorry}.

A caution on both counts. \ResProven{} out of \ResTotal{} is a high
proportion because a claim is not registered here until its argument is
written down; the register measures what has been recorded, not what
fraction of the problem is solved. Likewise \LeanTotal{} theorems is a
count of declarations, not of mathematical content --- see
\S\ref{app:repro}, which explains why a theorem count is not a coverage
measure and cites a case in this program's own archive where seventy-one
\texttt{sorry}-free modules turned out to prove tautologies.

% ---- begin gen_verification.tex ------------------------------------
% Generated by paper/master/generate_apparatus.py.  Do not edit.
\begin{longtable}{@{}p{0.44\linewidth} l l@{}}
\caption{Analytic results by recorded status and evidence level. \emph{Status} is what the argument establishes; \emph{evidence} is what artefact backs it. The two are deliberately independent: a result can be \texttt{proven} in status and \texttt{record-backed} in evidence, meaning the proof exists and the machine artefact does not.}\label{tab:results}\\
\toprule
Result & Status & Evidence \\
\midrule
\endfirsthead
\toprule Result & Status & Evidence \\ \midrule \endhead
\bottomrule
\endfoot
\texttt{\scriptsize ANISOTROPY\_\allowbreak{}THREE\_\allowbreak{}COORDINATE\_\allowbreak{}SHARP\_\allowbreak{}BOUND} & proven & analytic \\
\texttt{\scriptsize ARCHIVE\_\allowbreak{}CENTER\_\allowbreak{}DEFECT\_\allowbreak{}OBSTRUCTIONS} & proven & analytic \\
\texttt{\scriptsize ARCHIVE\_\allowbreak{}CONCENTRATION\_\allowbreak{}LOCALIZATION} & proven & analytic \\
\texttt{\scriptsize ARCHIVE\_\allowbreak{}FINITE\_\allowbreak{}RANGE\_\allowbreak{}INVERSE} & proven & analytic \\
\texttt{\scriptsize ARCHIVE\_\allowbreak{}GLOBAL\_\allowbreak{}PAIRING\_\allowbreak{}OBSTRUCTION} & proven & analytic \\
\texttt{\scriptsize ARCHIVE\_\allowbreak{}HORIZONTAL\_\allowbreak{}RESTRICTION\_\allowbreak{}BOUNDARY} & proven & analytic \\
\texttt{\scriptsize ARCHIVE\_\allowbreak{}LYAPUNOV\_\allowbreak{}PATCHING} & proven & analytic \\
\texttt{\scriptsize ARCHIVE\_\allowbreak{}REFLECTION\_\allowbreak{}PERMANENCE} & proven & analytic \\
\texttt{\scriptsize ARCHIVE\_\allowbreak{}SIGNED\_\allowbreak{}SU2\_\allowbreak{}DIFFUSION} & proven & analytic \\
\texttt{\scriptsize ARCHIVE\_\allowbreak{}SMOOTH\_\allowbreak{}PROFILES} & proven & analytic \\
\texttt{\scriptsize ARCHIVE\_\allowbreak{}SOURCE\_\allowbreak{}TILT\_\allowbreak{}AND\_\allowbreak{}ROOTED\_\allowbreak{}BOUNDS} & proven & analytic \\
\texttt{\scriptsize ASSEMBLED\_\allowbreak{}Q4\_\allowbreak{}COERCIVITY} & proven & analytic \\
\texttt{\scriptsize BALABAN\_\allowbreak{}MULTISCALE\_\allowbreak{}CAUCHY\_\allowbreak{}SUMMABILITY} & proven & analytic \\
\texttt{\scriptsize COMMON\_\allowbreak{}SPACE\_\allowbreak{}CONTRACTIVE\_\allowbreak{}DRIFT} & proven & analytic \\
\texttt{\scriptsize COMPACT\_\allowbreak{}ROTOR\_\allowbreak{}FAST\_\allowbreak{}GAP} & proven & analytic \\
\texttt{\scriptsize CONDITIONAL\_\allowbreak{}GRADIENT\_\allowbreak{}REPAIR} & proven & analytic \\
\texttt{\scriptsize CONDITIONAL\_\allowbreak{}QUANTUM\_\allowbreak{}PATH\_\allowbreak{}COVARIANCE} & proven & analytic \\
\texttt{\scriptsize CREATOR\_\allowbreak{}PARENT\_\allowbreak{}GAP} & proven & analytic \\
\texttt{\scriptsize CREATOR\_\allowbreak{}VELOCITY\_\allowbreak{}INVERSION} & proven & analytic \\
\texttt{\scriptsize DIMENSION\_\allowbreak{}FIVE\_\allowbreak{}OPERATOR\_\allowbreak{}IRRELEVANCE} & proven & analytic \\
\texttt{\scriptsize DYNAMIC\_\allowbreak{}CONDITIONAL\_\allowbreak{}CUBIC\_\allowbreak{}ENERGY} & proven & analytic \\
\texttt{\scriptsize EARLIER\_\allowbreak{}B6\_\allowbreak{}RETAINED\_\allowbreak{}KRYLOV\_\allowbreak{}CAMPAIGN} & conditional & numerical \\
\texttt{\scriptsize EARLIER\_\allowbreak{}CONDITIONAL\_\allowbreak{}SPECTRAL\_\allowbreak{}FLOOR} & proven & analytic \\
\texttt{\scriptsize EARLIER\_\allowbreak{}CORNER\_\allowbreak{}BLOCKING\_\allowbreak{}REFLECTION\_\allowbreak{}DEFECT} & proven & analytic \\
\texttt{\scriptsize EARLIER\_\allowbreak{}CUBIC\_\allowbreak{}QUOTIENT\_\allowbreak{}REGULARITY} & proven & analytic \\
\texttt{\scriptsize EARLIER\_\allowbreak{}D4\_\allowbreak{}DISJOINT\_\allowbreak{}STAPLE\_\allowbreak{}COORDINATES} & proven & analytic \\
\texttt{\scriptsize EARLIER\_\allowbreak{}DECORATED\_\allowbreak{}HISTORY\_\allowbreak{}MERGING} & proven & analytic \\
\texttt{\scriptsize EARLIER\_\allowbreak{}DETERMINANT\_\allowbreak{}PARITY\_\allowbreak{}REPAIR} & proven & analytic \\
\texttt{\scriptsize EARLIER\_\allowbreak{}EQUAL\_\allowbreak{}RAY\_\allowbreak{}IDENTIFICATION} & proven & analytic \\
\texttt{\scriptsize EARLIER\_\allowbreak{}FINITE\_\allowbreak{}ORDER\_\allowbreak{}SUPPORT} & proven & analytic \\
\texttt{\scriptsize EARLIER\_\allowbreak{}GRAM\_\allowbreak{}NULL\_\allowbreak{}OPERATOR\_\allowbreak{}QUOTIENT} & proven & analytic \\
\texttt{\scriptsize EARLIER\_\allowbreak{}ISOLATED\_\allowbreak{}POSITIVE\_\allowbreak{}MATRIX\_\allowbreak{}ATOM} & proven & analytic \\
\texttt{\scriptsize EARLIER\_\allowbreak{}JOINT\_\allowbreak{}RANK\_\allowbreak{}VOLUME\_\allowbreak{}SCALING} & proven & analytic \\
\texttt{\scriptsize EARLIER\_\allowbreak{}LOCALIZED\_\allowbreak{}FRAME\_\allowbreak{}OBSTRUCTION} & proven & analytic \\
\texttt{\scriptsize EARLIER\_\allowbreak{}MOVING\_\allowbreak{}ADJOINT\_\allowbreak{}FORCE\_\allowbreak{}DERIVATIVE} & proven & analytic \\
\texttt{\scriptsize EARLIER\_\allowbreak{}NESTED\_\allowbreak{}QUOTIENT\_\allowbreak{}REDUCTION} & proven & analytic \\
\texttt{\scriptsize EARLIER\_\allowbreak{}POSITIVE\_\allowbreak{}SECTOR\_\allowbreak{}PHASE\_\allowbreak{}ISOLATION} & proven & analytic \\
\texttt{\scriptsize EARLIER\_\allowbreak{}REFLECTION\_\allowbreak{}ADAPTED\_\allowbreak{}BLOCKING} & proven & analytic \\
\texttt{\scriptsize EARLIER\_\allowbreak{}SHELL\_\allowbreak{}SYMMETRY\_\allowbreak{}RITZ\_\allowbreak{}CONTROL} & proven & analytic \\
\texttt{\scriptsize EARLIER\_\allowbreak{}SPHERE\_\allowbreak{}EQUAL\_\allowbreak{}COUPLING\_\allowbreak{}COVARIANCE} & proven & analytic \\
\texttt{\scriptsize EARLIER\_\allowbreak{}SU3\_\allowbreak{}POSITIVITY\_\allowbreak{}BOUNDARIES} & proven & analytic \\
\texttt{\scriptsize EARLIER\_\allowbreak{}VSU\_\allowbreak{}BOUNDED\_\allowbreak{}DOMAIN\_\allowbreak{}WELLPOSEDNESS} & proven & analytic \\
\texttt{\scriptsize EXPONENTIAL\_\allowbreak{}SOURCE\_\allowbreak{}TOTALITY} & proven & analytic \\
\texttt{\scriptsize FESHBACH\_\allowbreak{}RESOLVENT\_\allowbreak{}COMPARISON} & proven & analytic \\
\texttt{\scriptsize FORM\_\allowbreak{}SCHUR\_\allowbreak{}SCALE\_\allowbreak{}COMPARISON} & proven & analytic \\
\texttt{\scriptsize FULL\_\allowbreak{}GRAPH\_\allowbreak{}SOURCE\_\allowbreak{}CONGRUENCE} & proven & analytic \\
\texttt{\scriptsize G17\_\allowbreak{}CORRECTED\_\allowbreak{}KP\_\allowbreak{}POLYMER\_\allowbreak{}STABILITY} & proven & analytic \\
\texttt{\scriptsize G9\_\allowbreak{}COMBINED\_\allowbreak{}MIXING\_\allowbreak{}SUPPORT} & proven & analytic \\
\texttt{\scriptsize G9\_\allowbreak{}HERMITIAN\_\allowbreak{}FOLDED\_\allowbreak{}SUPPORT} & proven & analytic \\
\texttt{\scriptsize G9\_\allowbreak{}LOCAL\_\allowbreak{}SIXTH\_\allowbreak{}NONCANCELLATION} & proven & analytic \\
\texttt{\scriptsize G9\_\allowbreak{}ONE\_\allowbreak{}FACE\_\allowbreak{}SIXTH} & proven & analytic \\
\texttt{\scriptsize G9\_\allowbreak{}SIXTH\_\allowbreak{}ORDER\_\allowbreak{}RUR\_\allowbreak{}ACTIVATION} & disputed & record-backed \\
\texttt{\scriptsize GAUSSIAN\_\allowbreak{}FULL\_\allowbreak{}FAST\_\allowbreak{}GREEN} & proven & analytic \\
\texttt{\scriptsize GAUSSIAN\_\allowbreak{}OS\_\allowbreak{}OBSERVABILITY} & proven & analytic \\
\texttt{\scriptsize GAUSSIAN\_\allowbreak{}PATH\_\allowbreak{}ENDPOINT\_\allowbreak{}BASELINE} & proven & analytic \\
\texttt{\scriptsize GAUSSIAN\_\allowbreak{}QUANTUM\_\allowbreak{}FAST\_\allowbreak{}SOURCES} & proven & analytic \\
\texttt{\scriptsize GAUSSIAN\_\allowbreak{}SCHUR\_\allowbreak{}MEMORY} & proven & analytic \\
\texttt{\scriptsize GROUND\_\allowbreak{}MARGINAL\_\allowbreak{}SCHUR\_\allowbreak{}SCORE} & proven & analytic \\
\texttt{\scriptsize HODGE\_\allowbreak{}FESHBACH\_\allowbreak{}SPLITTING} & proven & analytic \\
\texttt{\scriptsize HODGE\_\allowbreak{}INDUCED\_\allowbreak{}EIGHTH\_\allowbreak{}ORDER\_\allowbreak{}MOMENT} & proven & analytic \\
\texttt{\scriptsize HODGE\_\allowbreak{}LEAKAGE\_\allowbreak{}ZERO\_\allowbreak{}EXTENSION} & proven & analytic \\
\texttt{\scriptsize HODGE\_\allowbreak{}RANK\_\allowbreak{}ONE\_\allowbreak{}SECULAR\_\allowbreak{}CUBIC} & proven & analytic \\
\texttt{\scriptsize INTRINSIC\_\allowbreak{}SCORE\_\allowbreak{}CURRENT} & proven & analytic \\
\texttt{\scriptsize LITERAL\_\allowbreak{}ENDPOINT\_\allowbreak{}SPECTRAL\_\allowbreak{}COMPARISON} & proven & analytic \\
\texttt{\scriptsize LOCAL\_\allowbreak{}CLASS\_\allowbreak{}WICK\_\allowbreak{}COEFFICIENTS} & proven & analytic \\
\texttt{\scriptsize LOCAL\_\allowbreak{}CLASS\_\allowbreak{}WICK\_\allowbreak{}RECURRENCE} & proven & analytic \\
\texttt{\scriptsize LOCAL\_\allowbreak{}CLASS\_\allowbreak{}WICK\_\allowbreak{}SIGNS} & proven & analytic \\
\texttt{\scriptsize LOCAL\_\allowbreak{}GRADIENT\_\allowbreak{}EXCITATION\_\allowbreak{}SUPPORT} & proven & analytic \\
\texttt{\scriptsize MESOSCOPIC\_\allowbreak{}NORMALIZED\_\allowbreak{}SOURCE} & proven & analytic \\
\texttt{\scriptsize MOVING\_\allowbreak{}TIME\_\allowbreak{}MULTICHANNEL\_\allowbreak{}RATE} & proven & analytic \\
\texttt{\scriptsize MOVING\_\allowbreak{}TIME\_\allowbreak{}MULTICHANNEL\_\allowbreak{}SHARPNESS} & proven & analytic \\
\texttt{\scriptsize MOVING\_\allowbreak{}TIME\_\allowbreak{}SPECTRAL\_\allowbreak{}GAP} & proven & analytic \\
\texttt{\scriptsize ODD\_\allowbreak{}ORDER\_\allowbreak{}CENTRE\_\allowbreak{}PARITY} & proven & analytic \\
\texttt{\scriptsize ORDERED\_\allowbreak{}CONTOUR\_\allowbreak{}ACTIVITIES} & proven & analytic \\
\texttt{\scriptsize OS\_\allowbreak{}HISTORY\_\allowbreak{}BLOCK\_\allowbreak{}INTERTWINER} & proven & analytic \\
\texttt{\scriptsize OS\_\allowbreak{}KERNEL\_\allowbreak{}MOVING\_\allowbreak{}TIME\_\allowbreak{}GAP} & proven & analytic \\
\texttt{\scriptsize PREDICTABLE\_\allowbreak{}SOURCE\_\allowbreak{}DRIFT} & proven & analytic \\
\texttt{\scriptsize PROJECTIVE\_\allowbreak{}SOURCE\_\allowbreak{}MARTINGALE} & proven & analytic \\
\texttt{\scriptsize REVERSE\_\allowbreak{}MARTINGALE\_\allowbreak{}BLOCKING} & proven & analytic \\
\texttt{\scriptsize RICCATI\_\allowbreak{}RESIDUAL\_\allowbreak{}CERTIFICATE} & proven & analytic \\
\texttt{\scriptsize ROUGH\_\allowbreak{}GAUGE\_\allowbreak{}COERCIVITY\_\allowbreak{}LOWER\_\allowbreak{}BOUND} & proven & analytic \\
\texttt{\scriptsize RSM\_\allowbreak{}R\_\allowbreak{}CARRIER\_\allowbreak{}SYMBOL\_\allowbreak{}MASTER\_\allowbreak{}THEOREM} & proven & analytic \\
\texttt{\scriptsize SELECTED\_\allowbreak{}INVERSE\_\allowbreak{}WITHOUT\_\allowbreak{}DIAGONAL} & proven & analytic \\
\texttt{\scriptsize SHARP\_\allowbreak{}CUTOFF\_\allowbreak{}SOURCE\_\allowbreak{}GAP\_\allowbreak{}BUDGET} & proven & analytic \\
\texttt{\scriptsize SOURCE\_\allowbreak{}COMPLEMENT\_\allowbreak{}COMPLETE\_\allowbreak{}WINDOW} & proven & analytic \\
\texttt{\scriptsize SOURCE\_\allowbreak{}CONGRUENCE\_\allowbreak{}METRIC\_\allowbreak{}TRANSPORT} & proven & analytic \\
\texttt{\scriptsize SOURCE\_\allowbreak{}CURRENT\_\allowbreak{}BOUNDED\_\allowbreak{}W6} & proven & analytic \\
\texttt{\scriptsize SOURCE\_\allowbreak{}CURRENT\_\allowbreak{}VARIATIONAL\_\allowbreak{}ENERGY} & proven & analytic \\
\texttt{\scriptsize SQUARE\_\allowbreak{}FIRST\_\allowbreak{}CURRENT\_\allowbreak{}BUDGET} & proven & analytic \\
\texttt{\scriptsize SQUARE\_\allowbreak{}FIRST\_\allowbreak{}RESIDUAL\_\allowbreak{}CURRENT} & proven & analytic \\
\texttt{\scriptsize SQUARE\_\allowbreak{}SCALAR\_\allowbreak{}SOURCE\_\allowbreak{}ODD\_\allowbreak{}SCHUR\_\allowbreak{}VANISHING} & proven & analytic \\
\texttt{\scriptsize TETRAHEDRAL\_\allowbreak{}FLUX\_\allowbreak{}PROJECTION\_\allowbreak{}BOUNDARY} & proven & analytic \\
\texttt{\scriptsize TETRAHEDRAL\_\allowbreak{}S4\_\allowbreak{}COMMUTANT\_\allowbreak{}RESOLUTION} & proven & analytic \\
\texttt{\scriptsize TIER\_\allowbreak{}COLLAPSE\_\allowbreak{}ACTUAL\_\allowbreak{}H4\_\allowbreak{}SUPPORT} & proven & analytic \\
\texttt{\scriptsize TIER\_\allowbreak{}COLLAPSE\_\allowbreak{}IS\_\allowbreak{}R\_\allowbreak{}DEGREE} & falsified & analytic \\
\texttt{\scriptsize TRUE\_\allowbreak{}GROUND\_\allowbreak{}CENTER\_\allowbreak{}SCORE\_\allowbreak{}OBSTRUCTION} & proven & analytic \\
\texttt{\scriptsize UNIFORM\_\allowbreak{}MARKED\_\allowbreak{}CLUSTERING} & proven & analytic \\
\texttt{\scriptsize UNIVERSAL\_\allowbreak{}CELLULAR\_\allowbreak{}HODGE\_\allowbreak{}SPECTRUM} & proven & analytic \\
\texttt{\scriptsize W6\_\allowbreak{}ACTUAL\_\allowbreak{}CENTERED\_\allowbreak{}COMPLEMENT\_\allowbreak{}SUPPRESSION} & proven & analytic \\
\texttt{\scriptsize W6\_\allowbreak{}ANGULAR\_\allowbreak{}REFERENCE\_\allowbreak{}SCORE\_\allowbreak{}BUDGET} & proven & analytic \\
\texttt{\scriptsize W6\_\allowbreak{}R10\_\allowbreak{}CONDITIONAL\_\allowbreak{}DERIVATIVE} & proven & analytic \\
\texttt{\scriptsize W6\_\allowbreak{}R10\_\allowbreak{}H4\_\allowbreak{}FROM\_\allowbreak{}H0} & proven & analytic \\
\texttt{\scriptsize W6\_\allowbreak{}R10\_\allowbreak{}PROJECTION\_\allowbreak{}JET\_\allowbreak{}RECURRENCE} & proven & analytic \\
\texttt{\scriptsize W6\_\allowbreak{}UNIFORM\_\allowbreak{}AGMON\_\allowbreak{}COMPLEMENT\_\allowbreak{}GAP} & proven & analytic \\
\texttt{\scriptsize WILSON\_\allowbreak{}ACTIVITY\_\allowbreak{}EXTRACTION} & proven & analytic \\
\texttt{\scriptsize WILSON\_\allowbreak{}ACTUAL\_\allowbreak{}BLOCK\_\allowbreak{}FAST\_\allowbreak{}COMPLEMENT} & proven & analytic \\
\texttt{\scriptsize WILSON\_\allowbreak{}BLOCK\_\allowbreak{}SCORE\_\allowbreak{}OBSTRUCTION} & proven & analytic \\
\texttt{\scriptsize WILSON\_\allowbreak{}CARDINALITY\_\allowbreak{}CHART} & proven & analytic \\
\texttt{\scriptsize WILSON\_\allowbreak{}COEFFICIENT\_\allowbreak{}LOCALITY} & proven & analytic \\
\texttt{\scriptsize WILSON\_\allowbreak{}COMMON\_\allowbreak{}GAUSS\_\allowbreak{}LITERAL\_\allowbreak{}FAST\_\allowbreak{}COMPLEMENT} & proven & analytic \\
\texttt{\scriptsize WILSON\_\allowbreak{}COVARIANT\_\allowbreak{}BLOCK\_\allowbreak{}SOURCES} & proven & analytic \\
\texttt{\scriptsize WILSON\_\allowbreak{}CREATOR\_\allowbreak{}LIMIT} & proven & analytic \\
\texttt{\scriptsize WILSON\_\allowbreak{}CUBIC\_\allowbreak{}GROUND\_\allowbreak{}TRANSFER} & proven & analytic \\
\texttt{\scriptsize WILSON\_\allowbreak{}ENDPOINT\_\allowbreak{}COMPLETE\_\allowbreak{}CLUSTER} & proven & analytic \\
\texttt{\scriptsize WILSON\_\allowbreak{}ENDPOINT\_\allowbreak{}EQUATION} & proven & analytic \\
\texttt{\scriptsize WILSON\_\allowbreak{}FINITE\_\allowbreak{}CELL\_\allowbreak{}PHYSICAL\_\allowbreak{}GAP} & proven & analytic \\
\texttt{\scriptsize WILSON\_\allowbreak{}FIXED\_\allowbreak{}ORDER\_\allowbreak{}CHART} & proven & analytic \\
\texttt{\scriptsize WILSON\_\allowbreak{}FLAT\_\allowbreak{}BACKGROUND\_\allowbreak{}FAST\_\allowbreak{}SOURCES} & proven & analytic \\
\texttt{\scriptsize WILSON\_\allowbreak{}GLOBAL\_\allowbreak{}VERTICAL\_\allowbreak{}BARRIER} & proven & analytic \\
\texttt{\scriptsize WILSON\_\allowbreak{}GROUND\_\allowbreak{}BUNDLE\_\allowbreak{}RELATIVE\_\allowbreak{}FORM} & proven & analytic \\
\texttt{\scriptsize WILSON\_\allowbreak{}HARMONIC\_\allowbreak{}BOUNDARY} & proven & analytic \\
\texttt{\scriptsize WILSON\_\allowbreak{}HARMONIC\_\allowbreak{}CUBIC\_\allowbreak{}OBSTRUCTION} & proven & analytic \\
\texttt{\scriptsize WILSON\_\allowbreak{}INFINITE\_\allowbreak{}PHYSICAL\_\allowbreak{}BAND} & proven & analytic \\
\texttt{\scriptsize WILSON\_\allowbreak{}KINETIC\_\allowbreak{}WINDOW} & proven & analytic \\
\texttt{\scriptsize WILSON\_\allowbreak{}LITERAL\_\allowbreak{}COARSE\_\allowbreak{}FORM} & proven & analytic \\
\texttt{\scriptsize WILSON\_\allowbreak{}LITERAL\_\allowbreak{}FAST\_\allowbreak{}COMPLEMENT} & proven & analytic \\
\texttt{\scriptsize WILSON\_\allowbreak{}LOCALIZED\_\allowbreak{}TRUE\_\allowbreak{}GROUND\_\allowbreak{}SCORE} & proven & analytic \\
\texttt{\scriptsize WILSON\_\allowbreak{}PARENT\_\allowbreak{}GNS} & proven & analytic \\
\texttt{\scriptsize WILSON\_\allowbreak{}PHYSICAL\_\allowbreak{}SOURCE\_\allowbreak{}RANK\_\allowbreak{}REPAIR} & proven & analytic \\
\texttt{\scriptsize WILSON\_\allowbreak{}ROOTED\_\allowbreak{}CONTRACTION} & proven & analytic \\
\texttt{\scriptsize WILSON\_\allowbreak{}SAME\_\allowbreak{}WEIGHT\_\allowbreak{}OBSTRUCTION} & proven & analytic \\
\texttt{\scriptsize WILSON\_\allowbreak{}SECOND\_\allowbreak{}ORDER\_\allowbreak{}CHART} & proven & analytic \\
\texttt{\scriptsize WILSON\_\allowbreak{}STRIP\_\allowbreak{}BO\_\allowbreak{}FIRST\_\allowbreak{}TERM} & proven & analytic \\
\texttt{\scriptsize WILSON\_\allowbreak{}SYMMETRIC\_\allowbreak{}CREATORS} & proven & analytic \\
\texttt{\scriptsize WILSON\_\allowbreak{}THREE\_\allowbreak{}DIMENSIONAL\_\allowbreak{}HARMONIC\_\allowbreak{}BOUNDARY} & proven & analytic \\
\texttt{\scriptsize WILSON\_\allowbreak{}TRANSFER\_\allowbreak{}RESOLVENT} & proven & analytic \\
\texttt{\scriptsize WILSON\_\allowbreak{}TWO\_\allowbreak{}SQUARE\_\allowbreak{}PHYSICAL\_\allowbreak{}SHELLS} & proven & analytic \\
\texttt{\scriptsize WILSON\_\allowbreak{}TWO\_\allowbreak{}STRIP\_\allowbreak{}PHYSICAL\_\allowbreak{}SPLITTING} & proven & analytic \\
\texttt{\scriptsize WILSON\_\allowbreak{}UNIFORM\_\allowbreak{}FINITE\_\allowbreak{}SHELL} & proven & analytic \\
\texttt{\scriptsize WILSON\_\allowbreak{}VACUUM\_\allowbreak{}COMPRESSION} & proven & analytic \\
\texttt{\scriptsize WILSON\_\allowbreak{}VACUUM\_\allowbreak{}SPECTRAL\_\allowbreak{}FLOW} & proven & analytic \\
\texttt{\scriptsize WILSON\_\allowbreak{}VERTICAL\_\allowbreak{}FAST\_\allowbreak{}ENERGY} & proven & analytic \\
\texttt{\scriptsize WILSON\_\allowbreak{}WEIGHTED\_\allowbreak{}ACTIVITIES} & proven & analytic \\
\end{longtable}
% ---- end gen_verification.tex --------------------------------------
% ---- end A1_verification.tex ---------------------------------------

% ---- begin inlined bibliography ----------------------------------
\begin{thebibliography}{100}

\bibitem{AIZENMAN_1982}
Michael Aizenman.
\newblock Geometric analysis of $\varphi$$^4$ fields and ising models. parts i
  and ii.
\newblock {\em Communications in Mathematical Physics 86, 1-48}, 1982.

\bibitem{AT_2020_SU3}
Andreas Athenodorou and Michael Teper.
\newblock The glueball spectrum of su(3) gauge theory in 3+1 dimensions.
\newblock {\em JHEP 11 (2020) 172}, 2020.

\bibitem{AT_2021_SUN}
Andreas Athenodorou and Michael Teper.
\newblock Su(n) gauge theories in 3+1 dimensions: glueball spectrum, string
  tensions and topology.
\newblock {\em JHEP 12 (2021) 082}, 2021.

\bibitem{BV_1981_GEOMETRY}
O.~Babelon and C.~M. Viallet.
\newblock The riemannian geometry of the configuration space of gauge theories.
\newblock {\em Communications in Mathematical Physics 81(4), 515--525}, 1981.

\bibitem{BALABAN_1985}
T.~Balaban.
\newblock Ultraviolet stability of three-dimensional lattice pure gauge field
  theories.
\newblock {\em Communications in Mathematical Physics 102, 255-275}, 1985.

\bibitem{BALABAN_1987}
Tadeusz Ba\l{}aban.
\newblock Renormalization group approach to lattice gauge field theories. i.
  generation of effective actions in a small field approximation and a coupling
  constant renormalization in four dimensions.
\newblock {\em Communications in Mathematical Physics 109, 249-301}, 1987.

\bibitem{BBIJ_1984}
Tadeusz Balaban, John Imbrie, Arthur Jaffe, and David Brydges.
\newblock The mass gap for higgs models on a unit lattice.
\newblock {\em Annals of Physics 158, 281-319}, 1984.

\bibitem{BALAJI_2026}
S.~Balaji and et~al.
\newblock Quantum simulation of su(3) lattice yang-mills theory: one-cube
  benchmarks and circuit resource estimates, 2026.
\newblock arXiv preprint.

\bibitem{BRS_1976}
C~Becchi, A~Rouet, and R~Stora.
\newblock Renormalization of gauge theories.
\newblock {\em Annals of Physics 98, 287-321}, 1976.

\bibitem{BB_1983}
B.~Berg and A.~Billoire.
\newblock Glueball spectroscopy in 4d su(3) lattice gauge theory (i).
\newblock {\em Nucl. Phys. B221, 109-140}, 1983.

\bibitem{BORGA_2024}
Jacopo Borga, Sky Cao, and Jasper Shogren-Knaak.
\newblock Surface sums for lattice yang-mills in the large-n limit, 2024.
\newblock arXiv:2411.11676 [math.PR].

\bibitem{BFS_1979_I}
David Brydges, J\"urg Fr\"ohlich, and Erhard Seiler.
\newblock On the construction of quantized gauge fields. i. general results.
\newblock {\em Annals of Physics 121, 227-284}, 1979.

\bibitem{BF_1980}
David~C. Brydges and Paul Federbush.
\newblock Debye screening.
\newblock {\em Communications in Mathematical Physics 73, 197-246}, 1980.

\bibitem{BFS_1980_II}
David~C. Brydges, J\"urg Fr\"ohlich, and Erhard Seiler.
\newblock Construction of quantised gauge fields. ii. convergence of the
  lattice approximation.
\newblock {\em Communications in Mathematical Physics 71, 159-205}, 1980.

\bibitem{BFS_1981_III}
David~C. Brydges, J\"urg Fr\"ohlich, and Erhard Seiler.
\newblock On the construction of quantized gauge fields. iii. the
  two-dimensional abelian higgs model without cutoffs.
\newblock {\em Communications in Mathematical Physics 79, 353-399}, 1981.

\bibitem{CAO_2023}
Sky Cao, Minjae Park, and Scott Sheffield.
\newblock Random surfaces and lattice yang-mills.
\newblock {\em Commun. Amer. Math. Soc. 5 (2025) 774-896}, 2023.

\bibitem{CM_2003}
Jesse Carlsson and Bruce H.~J. McKellar.
\newblock Su(n) lattice gauge theory on a single cube, 2003.
\newblock arXiv preprint (unpublished).

\bibitem{CHEN_2006}
Y.~Chen, A.~Alexandru, S.J. Dong, T.~Draper, I.~Horvath, F.X. Lee, K.F. Liu,
  N.~Mathur, C.~Morningstar, M.~Peardon, S.~Tamhankar, B.L. Young, and J.B.
  Zhang.
\newblock Glueball spectrum and matrix elements on anisotropic lattices.
\newblock {\em Phys. Rev. D 73, 014516}, 2006.

\bibitem{CKS_2021}
Anthony Ciavarella, Natalie Klco, and Martin~J. Savage.
\newblock Trailhead for quantum simulation of su(3) yang-mills lattice gauge
  theory in the local multiplet basis.
\newblock {\em Phys. Rev. D 103, 094501}, 2021.

\bibitem{CB_2024}
Anthony~N. Ciavarella and Christian~W. Bauer.
\newblock Quantum simulation of su(3) lattice yang-mills theory at leading
  order in large-n\_c expansion.
\newblock {\em Phys. Rev. Lett. 133, 111901}, 2024.

\bibitem{CBB_2026}
Anthony~N. Ciavarella, Ivan~M. Burbano, and Christian~W. Bauer.
\newblock Efficient truncations of su(n\_c) lattice gauge theory for quantum
  simulation, 2026.
\newblock arXiv preprint.

\bibitem{BESIII_2026}
BESIII Collaboration.
\newblock Lightest 0-+ glueball as dominant constituent of x(2370), 2026.
\newblock arXiv:2607.20366 [hep-ex].

\bibitem{CS_2006}
Benoit Collins and Piotr Sniady.
\newblock Integration with respect to the haar measure on unitary, orthogonal
  and symplectic group.
\newblock {\em Comm. Math. Phys. 264, 773-795}, 2006.

\bibitem{CREUTZ_1980}
Michael Creutz.
\newblock Monte carlo study of quantized su(2) gauge theory.
\newblock {\em Physical Review D 21, 2308-2315}, 1980.

\bibitem{DRS_2021}
Zohreh Davoudi, Indrakshi Raychowdhury, and Andrew Shaw.
\newblock Search for efficient formulations for hamiltonian simulation of
  non-abelian lattice gauge theories.
\newblock {\em Phys. Rev. D 104, 074505}, 2021.

\bibitem{FP_1967}
L.D. Faddeev and V.N. Popov.
\newblock Feynman diagrams for the yang-mills field.
\newblock {\em Physics Letters B 25, 29-30}, 1967.

\bibitem{FLPS_2026}
Andrea Falzetti, Matteo Lombardi, Mauro~Lucio Papinutto, and Francesco
  Scardino.
\newblock Update on the computation of the quenched su(6) yang-mills lattice
  spectrum, 2026.
\newblock PoS LATTICE2025 (to appear); arXiv:2603.13138.

\bibitem{FO_1976}
Joel~S Feldman and Konrad Osterwalder.
\newblock The wightman axioms and the mass gap for weakly coupled ($\phi$4)3
  quantum field theories.
\newblock {\em Annals of Physics 97, 80-135}, 1976.

\bibitem{FEYNMAN_1981}
Richard~P. Feynman.
\newblock The qualitative behavior of yang-mills theory in 2 + 1 dimensions.
\newblock {\em Nuclear Physics B 188, 479-512}, 1981.

\bibitem{FOS_1983}
J.~Frohlich, K.~Osterwalder, and E.~Seiler.
\newblock On virtual representations of symmetric spaces and their analytic
  continuation.
\newblock {\em Annals of Mathematics 118, 461--489}, 1983.

\bibitem{FROHLICH_1982}
J\"urg Fr\"ohlich.
\newblock On the triviality of $\lambda$$\varphi$$^4$\_d theories and the
  approach to the critical point in d $\geq$ 4 dimensions.
\newblock {\em Nuclear Physics B 200, 281-296}, 1982.

\bibitem{GAWEDZKI_1985}
K.~Gaw\k{e}dzki and A.~Kupiainen.
\newblock Renormalization of a non-renormalizable quantum field theory.
\newblock {\em Nuclear Physics B 262, 33-48}, 1985.

\bibitem{GINIBRE_1970}
J.~Ginibre.
\newblock General formulation of griffiths' inequalities.
\newblock {\em Commun. Math. Phys. 16, 310-328}, 1970.

\bibitem{GJ_1968_I}
James Glimm and Arthur Jaffe.
\newblock A $\lambda$$\varphi$$^4$ quantum field theory without cutoffs. i.
\newblock {\em Physical Review 176, 1945--1951}, 1968.

\bibitem{GJ_1970_II}
James Glimm and Arthur Jaffe.
\newblock The $\lambda$($\varphi$$^4$)$_2$ quantum field theory without
  cutoffs. ii. the field operators and the approximate vacuum.
\newblock {\em Annals of Mathematics 91, 362--401}, 1970.

\bibitem{GJ_1970_III}
James Glimm and Arthur Jaffe.
\newblock The $\lambda$($\varphi$$^4$)$_2$ quantum field theory without
  cutoffs. iii. the physical vacuum.
\newblock {\em Acta Mathematica 125, 203-267}, 1970.

\bibitem{GJ_1972_IV}
James Glimm and Arthur Jaffe.
\newblock The $\lambda$$\varphi$$^4$$_2$ quantum field theory without cutoffs.
  iv. perturbations of the hamiltonian.
\newblock {\em Journal of Mathematical Physics 13, 1568-1584}, 1972.

\bibitem{GJ_1973_POSITIVITY}
James Glimm and Arthur Jaffe.
\newblock Positivity of the $\varphi$$^4$$_3$ hamiltonian.
\newblock {\em Fortschritte der Physik 21, 327-376}, 1973.

\bibitem{GJ_1974_SINGLE_PHASE}
James Glimm and Arthur Jaffe.
\newblock $\varphi$$^4$$_2$ quantum field model in the single-phase region:
  Differentiability of the mass and bounds on critical exponents.
\newblock {\em Physical Review D 10, 536-539}, 1974.

\bibitem{GJ_1985_SELECTED_I}
James Glimm and Arthur Jaffe.
\newblock Collected papers vol. 1: Quantum field theory and statistical
  mechanics: Expositions.
\newblock {\em Birkh\"auser Boston, ISBN 978-0-8176-3271-7}, 1985.

\bibitem{GJ_1985_SELECTED_II}
James Glimm and Arthur Jaffe.
\newblock Collected papers vol. 2: Constructive quantum field theory: Selected
  papers.
\newblock {\em Birkh\"auser Boston, ISBN 978-0-8176-3272-4}, 1985.

\bibitem{GJ_1987_QP}
James Glimm and Arthur Jaffe.
\newblock Quantum physics: A functional integral point of view (second
  edition).
\newblock {\em Springer New York, second edition}, 1987.

\bibitem{GJS_1974}
James Glimm, Arthur Jaffe, and Thomas Spencer.
\newblock The wightman axioms and particle structure in the p($\varphi$)$_2$
  quantum field model.
\newblock {\em Annals of Mathematics 100, 585--632}, 1974.

\bibitem{GJS_1976_I}
James Glimm, Arthur Jaffe, and Thomas Spencer.
\newblock A convergent expansion about mean field theory. i. the expansion.
\newblock {\em Annals of Physics 101, 610-630}, 1976.

\bibitem{GJS_1976_II}
James Glimm, Arthur Jaffe, and Thomas Spencer.
\newblock A convergent expansion about mean field theory. ii. convergence of
  the expansion.
\newblock {\em Annals of Physics 101, 631--669}, 1976.

\bibitem{GW_1973}
David~J. Gross and Frank Wilczek.
\newblock Ultraviolet behavior of non-abelian gauge theories.
\newblock {\em Physical Review Letters 30, 1343-1346}, 1973.

\bibitem{GRS_1975_I}
F.~Guerra, L.~Rosen, and B.~Simon.
\newblock The p($\varphi$)$_2$ euclidean quantum field theory as classical
  statistical mechanics. i.
\newblock {\em Annals of Mathematics 101, 111--189}, 1975.

\bibitem{GRS_1975_II}
F.~Guerra, L.~Rosen, and B.~Simon.
\newblock The p($\varphi$)$_2$ euclidean quantum field theory as classical
  statistical mechanics. ii.
\newblock {\em Annals of Mathematics 101, 191--259}, 1975.

\bibitem{GRS_1975_GAP}
Francesco Guerra, Lon Rosen, and Barry Simon.
\newblock Correlation inequalities and the mass gap in p($\varphi$)$_2$. iii.
  mass gap for a class of strongly coupled theories with nonzero external
  field.
\newblock {\em Communications in Mathematical Physics 41, 19-32}, 1975.

\bibitem{HAAG_1992}
Rudolf Haag.
\newblock Local quantum physics: Fields, particles, algebras.
\newblock {\em Springer, first edition}, 1992.

\bibitem{HAMER_1989}
C.~J. Hamer.
\newblock Hamiltonian strong coupling expansions for glueball masses in su(3).
\newblock {\em Phys. Lett. B 224, 339-342}, 1989.

\bibitem{HIP_1986}
C.~J. Hamer, A.~C. Irving, and T.~E. Preece.
\newblock Cluster expansion approach to non-abelian lattice gauge theory in
  (3+1)d (ii). su(3).
\newblock {\em Nucl. Phys. B270, 553}, 1986.

\bibitem{HSB_2000}
C.~J. Hamer, M.~Samaras, and R.~J. Bursill.
\newblock Green's function monte carlo study of su(3) lattice gauge theory in
  (3+1)d.
\newblock {\em Phys. Rev. D 62, 074506}, 2000.

\bibitem{HAZRA_2026}
Tamaghna Hazra.
\newblock New class of exactly flat topological bands - compact localised
  states protected by local graph topology.
\newblock {\em SciPost Physics (submission)}, 2026.

\bibitem{THOOFT_1974}
G.'t Hooft.
\newblock A planar diagram theory for strong interactions.
\newblock {\em Nuclear Physics B 72, 461-473}, 1974.

\bibitem{IMBRIE_1981_I}
John~Z. Imbrie.
\newblock Phase diagrams and cluster expansions for low temperature p(phi)\_2
  models. i. the phase diagram.
\newblock {\em Communications in Mathematical Physics 82, 261-304}, 1981.

\bibitem{IMBRIE_1981_II}
John~Z. Imbrie.
\newblock Phase diagrams and cluster expansions for low temperature p(phi)\_2
  models. ii. the schwinger functions.
\newblock {\em Communications in Mathematical Physics 82, 305-343}, 1981.

\bibitem{IH_1984}
A.~C. Irving and C.~J. Hamer.
\newblock Methods in hamiltonian lattice field theory (ii). linked-cluster
  expansions.
\newblock {\em Nucl. Phys. B230, 361}, 1984.

\bibitem{IZ_1980}
Claude Itzykson and Jean-Bernard Zuber.
\newblock Quantum field theory.
\newblock {\em McGraw-Hill, New York}, 1980.

\bibitem{JAFFE_2000_CQFT}
Arthur Jaffe.
\newblock Constructive quantum field theory.
\newblock {\em Mathematical Physics 2000, pp.111-127, Imperial College Press /
  World Scientific}, 2000.

\bibitem{JW_2006}
Arthur Jaffe and Edward Witten.
\newblock Quantum yang-mills theory.
\newblock {\em Clay Mathematics Institute, official 14-page problem statement},
  2006.

\bibitem{KRS_2023}
Saurabh~V. Kadam, Indrakshi Raychowdhury, and Jesse~R. Stryker.
\newblock Loop-string-hadron formulation of an su(3) gauge theory with
  dynamical quarks.
\newblock {\em Phys. Rev. D 107, 094513}, 2023.

\bibitem{KLEIN_ROSENBERGER_2020_WKB}
Markus Klein and Elke Rosenberger.
\newblock The tunneling effect for schr\"odinger operators on a vector bundle,
  2020.
\newblock arXiv preprint arXiv:2005.13852v1.

\bibitem{KPS_1981}
J.~B. Kogut, R.~B. Pearson, and J.~Shigemitsu.
\newblock The string tension, confinement and roughening in su(3) hamiltonian
  lattice gauge theory.
\newblock {\em Phys. Lett. B 98, 63}, 1981.

\bibitem{KSS_1976}
J.~B. Kogut, D.~K. Sinclair, and L.~Susskind.
\newblock A quantitative approach to low-energy quantum chromodynamics.
\newblock {\em Nucl. Phys. B114, 199-236}, 1976.

\bibitem{KS_1975}
J.~B. Kogut and L.~Susskind.
\newblock Hamiltonian formulation of wilson's lattice gauge theories.
\newblock {\em Phys. Rev. D 11, 395-408}, 1975.

\bibitem{KP_1986}
R.~Kotecky and D.~Preiss.
\newblock Cluster expansion for abstract polymer models.
\newblock {\em Communications in Mathematical Physics 103, 491-498}, 1986.

\bibitem{LEMOINE_2026}
Thibaut Lemoine.
\newblock Universal dualities for wilson loops in lattice yang-mills, 2026.
\newblock arXiv:2604.16252 [math-ph].

\bibitem{LLL_2006}
Mushtaq Loan, Xiang-Qian Luo, and Zhi-Huan Luo.
\newblock Monte carlo study of glueball masses in the hamiltonian limit of
  su(3) lattice gauge theory.
\newblock {\em Int. J. Mod. Phys. A 21, 2905}, 2006.

\bibitem{MRS_1993}
Jacques Magnen, Vincent Rivasseau, and Roland S\'en\'eor.
\newblock Construction of ym4 with an infrared cutoff.
\newblock {\em Communications in Mathematical Physics 155, 325-383}, 1993.

\bibitem{MS_1980_Y3}
Jacques Magnen and Roland S\'en\'eor.
\newblock Yukawa quantum field theory in three dimensions (y3).
\newblock {\em Annals of the New York Academy of Sciences 337, 13-43; Third
  International Conference on Collective Phenomena}, 1980.

\bibitem{MALDACENA_1998}
Juan~M. Maldacena.
\newblock The large n limit of superconformal field theories and supergravity.
\newblock {\em Advances in Theoretical and Mathematical Physics 2, 231-252},
  1998.

\bibitem{MR_1976_CRITICAL}
Oliver~A. McBryan and Jay Rosen.
\newblock Existence of the critical point in phi4 field theory.
\newblock {\em Communications in Mathematical Physics 51, 97-105}, 1976.

\bibitem{MMM_2019_CURVATURE}
Vincent Moncrief, Antonella Marini, and Rachel Maitra.
\newblock Orbit space curvature as a source of mass in quantum gauge theory,
  2019.
\newblock arXiv:1809.06318v2; Annals of Mathematical Sciences and Applications
  4(2).

\bibitem{MONDAL_2023_GEOMETRIC}
Puskar Mondal.
\newblock A geometric approach to the yang-mills mass gap, 2023.
\newblock arXiv:2301.06996v7.

\bibitem{MORNINGSTAR_2025}
Colin Morningstar.
\newblock Update on glueballs.
\newblock {\em PoS LATTICE2024 (2025) 004}, 2025.

\bibitem{MP_1999}
Colin~J. Morningstar and Mike Peardon.
\newblock The glueball spectrum from an anisotropic lattice study.
\newblock {\em Phys. Rev. D 60, 034509}, 1999.

\bibitem{MUNSTER_1981}
G.~Munster.
\newblock Strong coupling expansions for the mass gap in lattice gauge
  theories.
\newblock {\em Nucl. Phys. B190, 439}, 1981.

\bibitem{MUNSTER_1985_TM}
G.~Munster.
\newblock Effective transfer matrix for low-lying glueball states in lattice
  gauge theory.
\newblock {\em Nucl. Phys. B256, 67-76}, 1985.

\bibitem{NSY_2019}
B.~Nachtergaele, R.~Sims, and A.~Young.
\newblock Quasi-locality bounds for quantum lattice systems. part i.
  lieb-robinson bounds, quasi-local maps, and spectral flow automorphisms.
\newblock {\em J. Math. Phys. 60, 061101}, 2019.

\bibitem{NELSON_1973_MARKOFF}
Edward Nelson.
\newblock Quantum fields and markoff fields.
\newblock {\em Proceedings of Symposia in Pure Mathematics 23, 413-420,
  American Mathematical Society}, 1973.

\bibitem{OBZ_1985}
K.~H. O'Brien and J.-B. Zuber.
\newblock Strong coupling expansion of large-n qcd and surfaces.
\newblock {\em Nucl. Phys. B253, 621-634}, 1985.

\bibitem{ORAIFEARTAIGH_1997}
Lochlainn O'Raifeartaigh.
\newblock The dawning of gauge theory.
\newblock {\em Princeton University Press}, 1997.

\bibitem{OS_1973}
Konrad Osterwalder and Robert Schrader.
\newblock Axioms for euclidean green's functions.
\newblock {\em Communications in Mathematical Physics 31, 83-112}, 1973.

\bibitem{OS_1975}
Konrad Osterwalder and Robert Schrader.
\newblock Axioms for euclidean green's functions ii.
\newblock {\em Communications in Mathematical Physics 42, 281-305; appendix by
  Stephen Summers}, 1975.

\bibitem{OS_SEILER_1978}
Konrad Osterwalder and Erhard Seiler.
\newblock Gauge field theories on a lattice.
\newblock {\em Annals of Physics 110, 440-471}, 1978.

\bibitem{OS_SENEOR_1976}
Konrad Osterwalder and Roland S\'en\'eor.
\newblock The scattering matrix is non-trivial for weakly coupled p(phi)\_2
  models.
\newblock {\em Helvetica Physica Acta 49, 525-535}, 1976.

\bibitem{POLITZER_1973}
H.~David Politzer.
\newblock Reliable perturbative results for strong interactions?
\newblock {\em Physical Review Letters 30, 1346-1349}, 1973.

\bibitem{RIVASSEAU_1991}
Vincent Rivasseau.
\newblock From perturbative to constructive renormalization.
\newblock {\em Princeton University Press}, 1991.

\bibitem{SALAM_1968}
Abdus Salam.
\newblock Weak and electromagnetic interactions.
\newblock {\em Elementary Particle Theory, Proceedings of the Eighth Nobel
  Symposium, Lerum, Sweden, pp.367-377}, 1968.

\bibitem{SCHIERHOLZ_1988}
G.~Schierholz.
\newblock Glueball masses in su(3) lattice gauge theory.
\newblock {\em Nucl. Phys. B (Proc. Suppl.) 4, 11-20}, 1988.

\bibitem{SZH_1997}
D.~Schutte, W.-H. Zheng, and C.~J. Hamer.
\newblock The coupled cluster method in hamiltonian lattice field theory.
\newblock {\em Phys. Rev. D 55, 2974}, 1997.

\bibitem{SW_1994}
Nathan Seiberg and Edward Witten.
\newblock Electric-magnetic duality, monopole condensation, and confinement in
  n=2 supersymmetric yang-mills theory.
\newblock {\em Nuclear Physics B426, 19-52}, 1994.

\bibitem{SEILER_1975_Y2}
Erhard Seiler.
\newblock Schwinger functions for the yukawa model in two dimensions with
  space-time cutoff.
\newblock {\em Communications in Mathematical Physics 42, 163-182}, 1975.

\bibitem{SEO_1982}
K.~Seo.
\newblock Glueball mass estimate by strong coupling expansion in lattice gauge
  theories.
\newblock {\em Nucl. Phys. B209, 200-216}, 1982.

\bibitem{SIMON_1983_DISCRETE}
Barry Simon.
\newblock Some quantum operators with discrete spectrum but classically
  continuous spectrum.
\newblock {\em Annals of Physics 146, 209-220}, 1983.

\bibitem{SINGER_1981_ORBIT}
Isadore~M. Singer.
\newblock The geometry of the orbit space for non-abelian gauge theories.
\newblock {\em Physica Scripta 24, 817-820}, 1981.

\bibitem{STREATER_WIGHTMAN_1964}
Raymond~F. Streater and Arthur~S. Wightman.
\newblock Pct, spin and statistics, and all that.
\newblock {\em W. A. Benjamin, New York}, 1964.

\bibitem{SYMANZIK_1969}
Kurt Symanzik.
\newblock Euclidean quantum field theory.
\newblock {\em Local Quantum Theory, R. Jost (ed.), Academic Press, New York,
  pp.152-226}, 1969.

\bibitem{T_HOOFT_1979}
G.~'t~Hooft.
\newblock A property of electric and magnetic flux in nonabelian gauge
  theories.
\newblock {\em Nucl. Phys. B153, 141-160}, 1979.

\bibitem{UELTSCHI_2004_CLUSTER}
Daniel Ueltschi.
\newblock Cluster expansions and correlation functions.
\newblock {\em Moscow Mathematical Journal 4, 511-522}, 2004.

\bibitem{URZUA_2019}
Alejandro~R. Urzua, Iran Ramos-Prieto, Manuel Fernandez-Guasti, and Hector~M.
  Moya-Cessa.
\newblock Solution to the time-dependent coupled harmonic oscillators
  hamiltonian with arbitrary interactions.
\newblock {\em Quantum Reports 1, 82-90}, 2019.

\bibitem{WEINBERG_1967}
Steven Weinberg.
\newblock A model of leptons.
\newblock {\em Physical Review Letters19, 1264-1266}, 1967.

\bibitem{WILSON_1977_LECTURES}
Kenneth~G. Wilson.
\newblock Quarks and strings on a lattice.
\newblock {\em New Phenomena in Subnuclear Physics, Part A, The Subnuclear
  Series13, pp.69-142, Plenum}, 1977.

\bibitem{WITTEN_1994_FOUR_MANIFOLD}
Edward Witten.
\newblock Supersymmetric yang-mills theory on a four-manifold.
\newblock {\em Journal of Mathematical Physics35, 5101-5135}, 1994.

\bibitem{YANG_MILLS_1954}
Chen~Ning Yang and Robert~L. Mills.
\newblock Conservation of isotopic spin and isotopic gauge invariance.
\newblock {\em Physical Review96, 191-195}, 1954.

\bibitem{YAROTSKY_2004_GS}
D.~A. Yarotsky.
\newblock Ground states in relatively bounded quantum perturbations of
  classical lattice systems, 2004.
\newblock arXiv preprint.

\bibitem{YAROTSKY_2004_QP}
D.~A. Yarotsky.
\newblock Quasi-particles in weak perturbations of non-interacting quantum
  lattice systems, 2004.
\newblock arXiv preprint.

\end{thebibliography}
% ---- end inlined bibliography ------------------------------------
\end{document}
